% Options for packages loaded elsewhere
\PassOptionsToPackage{unicode}{hyperref}
\PassOptionsToPackage{hyphens}{url}
%
\documentclass[
]{article}
\usepackage{amsmath,amssymb}
\usepackage{iftex}
\ifPDFTeX
  \usepackage[T1]{fontenc}
  \usepackage[utf8]{inputenc}
  \usepackage{textcomp} % provide euro and other symbols
\else % if luatex or xetex
  \usepackage{unicode-math} % this also loads fontspec
  \defaultfontfeatures{Scale=MatchLowercase}
  \defaultfontfeatures[\rmfamily]{Ligatures=TeX,Scale=1}
\fi
\usepackage{lmodern}
\ifPDFTeX\else
  % xetex/luatex font selection
\fi
% Use upquote if available, for straight quotes in verbatim environments
\IfFileExists{upquote.sty}{\usepackage{upquote}}{}
\IfFileExists{microtype.sty}{% use microtype if available
  \usepackage[]{microtype}
  \UseMicrotypeSet[protrusion]{basicmath} % disable protrusion for tt fonts
}{}
\makeatletter
\@ifundefined{KOMAClassName}{% if non-KOMA class
  \IfFileExists{parskip.sty}{%
    \usepackage{parskip}
  }{% else
    \setlength{\parindent}{0pt}
    \setlength{\parskip}{6pt plus 2pt minus 1pt}}
}{% if KOMA class
  \KOMAoptions{parskip=half}}
\makeatother
\usepackage{xcolor}
\setlength{\emergencystretch}{3em} % prevent overfull lines
\providecommand{\tightlist}{%
  \setlength{\itemsep}{0pt}\setlength{\parskip}{0pt}}
\setcounter{secnumdepth}{-\maxdimen} % remove section numbering
% unicode-fix.tex — render raw-Unicode symbols that the Latin Modern *text* font
% lacks, using Computer Modern math/text equivalents. This keeps the established
% Computer Modern look while filling the glyphs xelatex would otherwise drop.
% Included via: pandoc ... --include-in-header=unicode-fix.tex
\usepackage{newunicodechar}
\usepackage{pifont}
\usepackage{graphicx} % unicode-math's smallsetminus uses rotatebox
% Use the Computer Modern-compatible set-difference construction; the large
% Unicode glyph selected by TeX Live 2022 is absent from Latin Modern Math.
\AtBeginDocument{\let\setminus\smallsetminus}
% --- Greek letters -> math (Latin Modern Math; consistent with CM body text) ---
\newunicodechar{α}{\ensuremath{\alpha}}
\newunicodechar{β}{\ensuremath{\beta}}
\newunicodechar{γ}{\ensuremath{\gamma}}
\newunicodechar{δ}{\ensuremath{\delta}}
\newunicodechar{ε}{\ensuremath{\varepsilon}}
\newunicodechar{θ}{\ensuremath{\theta}}
\newunicodechar{μ}{\ensuremath{\mu}}
\newunicodechar{ν}{\ensuremath{\nu}}
\newunicodechar{ρ}{\ensuremath{\rho}}
\newunicodechar{σ}{\ensuremath{\sigma}}
\newunicodechar{τ}{\ensuremath{\tau}}
\newunicodechar{φ}{\ensuremath{\varphi}}
\newunicodechar{χ}{\ensuremath{\chi}}
\newunicodechar{ψ}{\ensuremath{\psi}}
\newunicodechar{η}{\ensuremath{\eta}}
\newunicodechar{κ}{\ensuremath{\kappa}}
\newunicodechar{λ}{\ensuremath{\lambda}}
\newunicodechar{π}{\ensuremath{\pi}}
\newunicodechar{ω}{\ensuremath{\omega}}
\newunicodechar{Ḣ}{\ensuremath{\dot H}}
\newunicodechar{≪}{\ensuremath{\ll}}
\newunicodechar{≫}{\ensuremath{\gg}}
\newunicodechar{⟺}{\ensuremath{\Longleftrightarrow}}
\newunicodechar{⟹}{\ensuremath{\Longrightarrow}}
\newunicodechar{⇔}{\ensuremath{\Leftrightarrow}}
\newunicodechar{⟨}{\ensuremath{\langle}}
\newunicodechar{⟩}{\ensuremath{\rangle}}
\newunicodechar{⁰}{\textsuperscript{0}}
\newunicodechar{¹}{\textsuperscript{1}}
\newunicodechar{²}{\textsuperscript{2}}
\newunicodechar{³}{\textsuperscript{3}}
\newunicodechar{⁴}{\textsuperscript{4}}
\newunicodechar{⁵}{\textsuperscript{5}}
\newunicodechar{⁶}{\textsuperscript{6}}
\newunicodechar{⁷}{\textsuperscript{7}}
\newunicodechar{⁸}{\textsuperscript{8}}
\newunicodechar{⁹}{\textsuperscript{9}}
% --- Relations / operators -> math ---
\newunicodechar{↔}{\ensuremath{\leftrightarrow}}
\newunicodechar{⇐}{\ensuremath{\Leftarrow}}
\newunicodechar{⇒}{\ensuremath{\Rightarrow}}
\newunicodechar{∂}{\ensuremath{\partial}}
\newunicodechar{∅}{\ensuremath{\varnothing}}
\newunicodechar{∈}{\ensuremath{\in}}
\newunicodechar{∼}{\ensuremath{\sim}}
\newunicodechar{≈}{\ensuremath{\approx}}
\newunicodechar{≠}{\ensuremath{\neq}}
\newunicodechar{≤}{\ensuremath{\leq}}
\newunicodechar{≥}{\ensuremath{\geq}}
\newunicodechar{⊂}{\ensuremath{\subset}}
\newunicodechar{⊃}{\ensuremath{\supset}}
% --- Special ---
\newunicodechar{ℏ}{\ensuremath{\hbar}}
\newunicodechar{ℤ}{\ensuremath{\mathbb{Z}}}
% --- Superscript modifier letters (footnote-style markers) -> text superscript ---
\newunicodechar{ʰ}{\textsuperscript{h}}
\newunicodechar{ʲ}{\textsuperscript{j}}
\newunicodechar{ᵃ}{\textsuperscript{a}}
\newunicodechar{ᵇ}{\textsuperscript{b}}
\newunicodechar{ᵈ}{\textsuperscript{d}}
\newunicodechar{ᵉ}{\textsuperscript{e}}
\newunicodechar{ᵍ}{\textsuperscript{g}}
\newunicodechar{ᶠ}{\textsuperscript{f}}
\newunicodechar{⁻}{\textsuperscript{\ensuremath{-}}}
% --- Subscript digits -> text subscript ---
\newunicodechar{₀}{\textsubscript{0}}
\newunicodechar{₁}{\textsubscript{1}}
\newunicodechar{₂}{\textsubscript{2}}
\newunicodechar{₃}{\textsubscript{3}}
\newunicodechar{₄}{\textsubscript{4}}
\newunicodechar{₈}{\textsubscript{8}}
% --- Checkmarks ---
\newunicodechar{✓}{\ding{51}}
\newunicodechar{✗}{\ding{55}}

% --- added 2026-08-02 after an independent glyph audit ---
\newunicodechar{ℓ}{\ensuremath{\ell}}
\newunicodechar{⊕}{\ensuremath{\oplus}}
\newunicodechar{⊗}{\ensuremath{\otimes}}
\newunicodechar{∠}{\ensuremath{\angle}}
\newunicodechar{≡}{\ensuremath{\equiv}}
\newunicodechar{∝}{\ensuremath{\propto}}
\newunicodechar{ℝ}{\ensuremath{\mathbb{R}}}
\newunicodechar{₅}{\ensuremath{_5}}
\newunicodechar{₋}{\ensuremath{_-}}
\newunicodechar{⁺}{\ensuremath{^+}}
\newunicodechar{⁽}{\ensuremath{^(}}
\newunicodechar{⁾}{\ensuremath{^)}}
% U+239B..U+23A0: multi-line parenthesis assembly pieces emitted by some
% converters. There is no sane inline equivalent; map to ordinary delimiters so
% the notation survives rather than vanishing silently.
\newunicodechar{⎛}{\ensuremath{\big(}}
\newunicodechar{⎜}{\ensuremath{\big\vert}}
\newunicodechar{⎝}{\ensuremath{\big(}}
\newunicodechar{⎞}{\ensuremath{\big)}}
\newunicodechar{⎟}{\ensuremath{\big\vert}}
\newunicodechar{⎠}{\ensuremath{\big)}}
\ifLuaTeX
  \usepackage{selnolig}  % disable illegal ligatures
\fi
\IfFileExists{bookmark.sty}{\usepackage{bookmark}}{\usepackage{hyperref}}
\IfFileExists{xurl.sty}{\usepackage{xurl}}{} % add URL line breaks if available
\urlstyle{same}
\hypersetup{
  hidelinks,
  pdfcreator={LaTeX via pandoc}}

\author{}
\date{}

\begin{document}

\hypertarget{the-substratum-construction-reconstruction-the-substratum-gauge-group-and-the-qm-gr-synthesis}{%
\section{The Substratum Construction: Reconstruction, the Substratum
Gauge Group, and the QM-GR
Synthesis}\label{the-substratum-construction-reconstruction-the-substratum-gauge-group-and-the-qm-gr-synthesis}}

\textbf{Observer geometry and physical carrier.} SM §4.1 distinguishes
two claims. First, finite equivariant projection is classical and exact:
cubic symmetry is inherited by the projected Markov part and every
memory kernel, forbidding quadratic anisotropy in the scalar symbol ---
so whatever nondegenerate quadratic structure survives is isotropic ---
while allowing cubic anisotropy at quartic order. The particular
normalized nearest-neighbor operator \(A/(2d)\) additionally requires an
\textbf{observer-level} center-free kernel; microscopic center
independence alone does not imply that after correlated hidden variables
are marginalized. Second, the six signed simple-cubic links and their
exact \(3\oplus2\oplus1\) representation are unconditional geometry, but
identifying that link space with the complete physical carrier is
\textbf{H-link}; the single-copy clause gives \(K=6\), and the custodial
condensate-stabilizer reading is \textbf{H-cust}. Accordingly, any
downstream statement here that consumes \(K=6\), \(N_f=6\), the physical
SM gauge carrier, or a fixed three-sector count inherits those named
bridge conditions unless it is explicitly referring only to the six
geometric links.

\textbf{Author:} Alex Maybaum\\
\textbf{Date:} April 2026\\
\textbf{Status:} DRAFT PRE-PRINT\\
\textbf{Classification:} Theoretical Physics / Foundations

\begin{center}\rule{0.5\linewidth}{0.5pt}\end{center}

\hypertarget{abstract}{%
\subsection{Abstract}\label{abstract}}

Quantum mechanics and general relativity are conventionally treated as
two fundamental theories awaiting unification. This paper develops the
alternative that they are two projections of a single finite
deterministic construction, joined by a third projection that produces
the arrow of time. The underlying object is the bijection
\((S,\varphi)\) on a finite configuration space. Its visible-sector
reduction has the exact finite stochastic/quantum-representation
equivalence established in {[}Main{]}: finite observable laws, finite
reversible realizations, and fixed-basis unitary/Born representations
coincide, while genuine history readback forces hidden predictive
memory. Its spatial coupling graph supplies finite causal cones;
equality with the complete standard local quantum instrument theory is
the narrower operational-lifting statement explicitly scoped in
{[}Main{]}. The boundary-thermodynamic projection yields the
gravitational description developed in {[}GR{]}, and the lattice
dynamics yields the Standard-Model structure developed in {[}SM{]}.

Theorem 23 collects conditional local lattice/gauge outputs of a
selected cubic branch under its named physical and structural premises.
It does not select that branch uniquely from E1--E7 and A1--A6; a
Bell-compatible common completion additionally requires H-Bell. Theorem
24 establishes exact hidden-only coherent-channel invariance, broader
rooted-population invariance and independent decoupled enlargement,
together with separately scoped alphabet and graph comparisons. Pair C
refutes exhaustiveness of the restricted coherent-channel generators,
while equivalence across all admissible partitions remains open. These
results motivate a common-substratum interpretation of the quantum,
lattice-gauge and horizon constructions. They do not establish a
complete observational quotient or a complete projection of its
transformations onto the Standard Model gauge group.

\hypertarget{introduction}{%
\subsection{1. Introduction}\label{introduction}}

The standard physical picture of the world contains two fundamental
theories. Quantum mechanics describes the microscopic regime in terms of
unitary evolution on a Hilbert space, superposition, and Born-rule
probabilities. General relativity describes the gravitational regime in
terms of a deterministic, classical, dynamical metric obeying Einstein's
equations. The two are mathematically and conceptually incompatible.
Quantum mechanics has no preferred temporal foliation and no
observer-independent state; general relativity has both. Quantum
mechanics treats measurement as a primitive operation; general
relativity treats it as a process within the dynamics. Quantum mechanics
is linear; general relativity is not. Every attempt to unify them --- to
quantize gravity, to geometrize quantum mechanics, to merge the two into
a deeper underlying theory --- has produced either inconsistencies,
ambiguities, or empirically unfalsifiable conjectures. The unification
problem has resisted solution for nearly a century, and the most active
current programs (loop quantum gravity, string theory, causal dynamical
triangulations, asymptotic safety) have produced neither a quantitative
empirical confirmation nor a clean conceptual resolution.

This paper argues that the QM-GR incompatibility is not a technical
problem to be solved by finding the right mathematical framework. It is
a category error: an attempt to merge two descriptions that turn out to
be projections of the same underlying object viewed from different
perspectives, rather than two competing accounts of the same regime. The
framework developed in this four-paper synthesis to which this paper
belongs --- {[}Main{]}, {[}SM{]}, {[}GR{]}, and the present paper ---
provides the underlying object and the projection map. The underlying
object is a finite deterministic bijection \((S, \varphi)\) on a
configuration space partitioned into a visible sector accessible to an
embedded observer and a hidden sector beyond the observer's causal
reach. The projection from the underlying object to the observer's
reduced description is the trace-out over the hidden sector. {[}Main{]}
proves the exact finite observable-law equivalence
\(S\iff D\iff Q_{\mathrm{fb}}\): finite stochastic laws, finite
reversible realizations, and fixed-basis unitary/Born representations
coincide. The representation itself is universal; the memory-bearing OI
sector is singled out instead by history readback, which forces hidden
predictive memory. For spatial systems the same construction supplies
graph causal cones for finite-range couplings. What remains conditional
on the quantum side is the strongest operational reading --- one common
standard quantum instrument/composite theory for the whole coherent
intervention family. Bell is separate: the response-complete
deterministic layer keeps measurement independence and reproduces
quantum Bell violation only on the ontically parameter-dependent branch
of {[}Main §3.3{]}, with operational no-signaling. The projection from
the same underlying object to general relativity is the thermodynamic
limit at the partition boundary, formalized in {[}GR{]} as the
cosmological-horizon application: applied to the cosmological horizon as
a causal partition, the framework relates the action scale to the
discreteness scale through the calibration
\(\hbar = c^3 \epsilon^2 / (4G\sigma)\), which fixes only
\(\epsilon/\sqrt{\sigma} = 2\,l_p\) for an entropy \(\sigma\) per
boundary mode, and recovers the Bekenstein-Hawking entropy with the
\(1/4\) coefficient as a consistency relation of that calibration, both
conditional on H-slope {[}GR §3.2{]}, and proposes a dissolution of the
cosmological constant problem conditional on the covariant coupling step
{[}GR §6.2{]}. The Standard Model emerges from the same bijection on a
cubic lattice with the wave equation as substratum dynamics: the cubic
lattice is selected by the observed gauge group (the \(\mathrm{SU}(3)\)
factor requires the cubic crystal system), and --- given that lattice
--- three generations (physical reading conditional on H-spin', {[}SM
§4.7{]}; chirality on H-χ', {[}SM §4.8{]}), the Higgs as composite, and
the \((3,2,1)\) gauge structure follow as forced consequences {[}SM{]};
strong CP is open there (H-top, H-det). The framework's hierarchical
structural realism and its relationship to other unification programs
are developed in the fifth core paper {[}Structure{]}.

The present paper compares the visible-law construction of {[}Main{]},
the selected lattice realization of {[}SM{]} and the horizon
construction of {[}GR{]}. Theorem 23 (§3) states conditional outputs of
the nearest-neighbor cubic wave-rule branch under H-link, the octahedral
representation, amplitude-scale gauge and the separate physical
hypotheses used by {[}SM{]} and {[}GR{]}. The converse from the
empirical inputs to a unique local branch is unproved, even within
A1--A6 without further selection assumptions. H-Bell separately
conditions a common Bell-compatible completion. Section 4 distinguishes
exact coherent-channel invariances from population-only invariance and
structural alphabet/graph comparisons. The restricted coherent-channel
generators are not exhaustive; the all-partitions classification and a
complete Standard Model gauge-image map remain open. The proposed
synthesis shares a model and observation machinery while retaining each
construction's additional premises.

The synthesis proposed in §5 treats the quantum-representable visible
law, the selected cubic gauge construction and horizon thermodynamics as
descriptions associated with a common candidate substratum. {[}Main{]}
supplies the exact finite observable-law representation. The physical
Standard Model interpretation retains the hypotheses of {[}SM{]}, while
the action-scale and gravitational conclusions retain H-slope, the
horizon/frame premises and the coupling conditions of {[}GR{]}. Sharing
a candidate object does not prove that these descriptions constitute a
unique common physical realization. Theorem 23 does not select that
realization from observations, and Theorem 24 does not furnish a
complete gauge quotient. The horizon-growth account of the observer's
arrows (§5.4) likewise retains its cosmological and observer-admission
conditions.

The paper is organized as follows. §1.1 states the framework's primary
metaphysical commitment and its relation to explanatory and empirical
claims. §2 develops the physical interpretation of \((S, \varphi)\) as a
finite lossless memory with partial read-write access, and addresses the
substrate objection (what is the memory made of?) by distinguishing
demonstrated observational redundancies from the open question of deeper
physical structure (§4). §3 states and proves the reconstruction theorem
(Theorem 23). §4 proves selected fixed-interface invariances (Theorem
24), gives the restricted block-word rigidity result (Proposition 24.2),
and examines the proposed gauge hierarchy with its open completeness
question. §5 makes the synthesis claim explicit, citing {[}Main{]},
{[}SM{]}, and {[}GR{]} for the empirical content. §6 discusses
structural realism, the ontological hierarchy, and the dissolution of
the measurement problem. §7 concludes. The paper is short by design: the
technical content fits in two sections (§§3--4) and the rest is the
structural and interpretive context that shows why those two sections
matter.

\hypertarget{the-metaphysical-commitment}{%
\subsubsection{1.1 The metaphysical
commitment}\label{the-metaphysical-commitment}}

The framework's primary metaphysical proposal is ontic structural
realism about the selected local model and the structures preserved by
demonstrated observational redundancies. Theorem 23 supplies conditional
outputs of that branch; it does not classify all microscopic
representatives compatible with observed physics. Section 4 identifies
invariances only for specified observational carriers. H-Bell is a
separate compatibility condition for a common Bell-violating completion,
and neither Theorem 23 nor Proposition 24.2 establishes its uniqueness.
Commitment to one complete microscopic ontology therefore requires more
than these results.

Stated at its most distilled, the framework's foundational commitment
--- what {[}Main §1.2{]} gives as ``observation occurs'' --- is itself
structured: it is two axioms. The first, that tokened differentiation
occurs, is indubitable and conceptually primitive. The second, that
differentiation recurs --- across more than one domain, minimally
configuration and time --- is a substantive posit: not indubitable, not
derivable from the first (a no-go result establishes this), though
presupposing it. The commitment is single in the sense that the two
axioms invoke one principle; it is, in count, two, and the second is the
framework's one non-evidential foundational posit. Part I of the
companion paper \emph{Physics Modulo Gauge} itemizes what the first
axiom forces as lemmas and what the second adds. (One item of that
itemization is a reading of the primitive, adopted here explicitly: the
observer's embeddedness --- the proper-part decomposition
\(V \subsetneq S\) --- belongs to the first axiom, because ``tokened
differentiation'' is read perspectivally, as a registering from a locus
against a remainder. On a thin reading it would instead be a third
axiom; the framework adopts the perspectival reading, so the count is
two. The choice is presentational, not empirical, and changes nothing
downstream of \((S, \varphi, V)\).)

The structural-realism commitment operates at multiple levels. Within
OI's proposed structural class, the realism about
\([(S, \varphi)] / \mathcal{G}_{\text{sub}}\) depends on a specified
observable equivalence and an unresolved classification of the
transformations preserving it. Across structural classes --- at the
level of partial-trace observational features that any embedded-observer
system must respect --- broader structural realism applies, with the
framework's universality-class equivalence ({[}Structure §9{]}) being
the appropriate gauge equivalence. The full multi-level structure is
articulated in {[}Structure{]}, which canonicalizes both the observation
hierarchy and the gauge hierarchy and their orthogonality, and develops
the cross-framework analysis through universality classes of embedded
observers. The substratum-class commitment of this paper is one level of
a broader hierarchical structural realism: realism at multiple
combinations of observation-hierarchy depth and gauge-hierarchy breadth,
with the framework's empirical content concentrated at the deepest
observation level (Level D in {[}Structure §2{]}) and the most
class-specific gauge level (Level G3 in {[}Structure §2{]}). The
articulation here focuses on this concentrated level because that is
where the reconstruction theorem and substratum gauge group operate; the
broader hierarchical content is developed in {[}Structure{]}.

The structural-realist proposal has explanatory and empirical aspects.
The exact finite-law representation in {[}Main §3{]} explains how a
selected reversible model produces a quantum representation of its
accessible law. Within the selected physical branch, the representation
calculations and their named hypotheses constrain lattice/gauge outputs
and the quantities classified in {[}SM §7{]} and {[}GR §7{]}. Their
predictive or retrodictive status is that of the individual derivation
chains. Neither explanatory economy nor agreement with those quantities
supplies the missing converse from observations to a unique substratum.
The structural, explanatory and empirical questions must therefore be
assessed separately.

\begin{center}\rule{0.5\linewidth}{0.5pt}\end{center}

\hypertarget{the-physical-interpretation-of-s-ux3c6}{%
\subsection{2. The Physical Interpretation of (S,
φ)}\label{the-physical-interpretation-of-s-ux3c6}}

\textbf{Storage and memory.} S is the set of all distinguishable states:
\emph{finite capacity}. φ is a bijection: \emph{perfect memory} ---
information is never created or destroyed. Together, \((S,\varphi)\) is
a finite lossless memory. The partition \(V\) defines the observer's
\emph{access} --- both read and write. The observer reads the visible
sector and, through the coupling \(H_{\text{int}}\), visible operations
write correlations into the hidden sector. The slow-bath form of (C2) is
one mechanism that preserves those writes until readback. C4 then makes
the visible law genuinely history-sensitive; on a fixed finite recurrent
representative any non-permutation rooted map, which that readback
supplies, forces indivisibility somewhere in the recurrence cycle. The
resulting finite observable law has an exact fixed-basis unitary/Born
representation by the foundational equivalence theorem. The unresolved
step is not the existence or quadratic form of that representation, but
its simultaneous extension to the full coherent local
instrument/composite algebra. Interference in the mechanism picture is
write-then-read: information deposited in the hidden sector during one
transition is retrieved later and changes the visible transition
statistics.

The \(10^{122}\) CC discrepancy is the compression ratio between total
storage and readable storage. The dark sector is the gravitational
effect of the unreadable storage. The Bekenstein-Hawking entropy is the
storage capacity of the partition boundary. The action scale ℏ is the
conversion factor between storage geometry and read-write statistics.

\textbf{Remark (the domain of the local lattice dynamics).} The wave
equation of {[}SM §4.1{]} --- the unique second-order reversible
nearest-neighbor dynamics compatible with center independence, isotropy,
and linearity --- governs the full \textbf{local lattice sector} used by
the SM/GR construction, not the visible degrees of freedom alone.
Visible and hidden lattice degrees of freedom in that sector run the
same wave equation; the partition \(V\) imposes observer access without
changing the local update. This is enough for (i) the trace-out of
{[}Main{]} over lattice-hidden degrees of freedom, (ii) the nested local
trace-out of {[}GR §8.4{]}, and (iii) the link-carrier construction of
{[}SM §4.7.1.2{]}. It is not enough for the Bell-violating deterministic
completion under measurement independence: a nearest-neighbor update,
whose cone advances one edge per update, together with interventions
each reading and writing only its own wing and readouts outside the
remote setting's graph cone, falls under {[}Main §3.3{]}'s Bell ceiling
and gives \(|S_{\rm CHSH}|\le2\). The one-object synthesis therefore
additionally requires H-Bell (§3.3): the state-dependent graph must
realize Bell-relevant preparation-indexed parameter dependence with
operational no-signaling while preserving the metric/Ollivier--Ricci
continuum limit required by the GR projection.

\textbf{The substrate objection.} ``What is the memory made of?'' A
specified \((S,\varphi)\), observer partition and preparation/readout
rule determine a model's accessible law. Theorem 23 gives conditional
local outputs of a selected realization, not a proof that this model is
a complete description of reality or is uniquely recovered from
observations. If two interpretations assign exactly the same accessible
data to that model, any measurement procedure depending only on those
data assigns them the same result, by the factoring argument {[}GR
§6.3{]}. This conditional indistinguishability does not establish that
every proposed microscopic difference is gauge. Questions about a deeper
physical realization remain distinct from interpretive readings that
stipulate identical data.

\textbf{Relation to computation.} A Turing machine has a tape (storage),
a head (partial read/write access), and a transition function (update
rule). The correspondence is suggestive: S is the tape, V is the head's
access window, φ is the transition function. But three differences are
physically significant. First, a Turing machine's tape is potentially
infinite; S is finite --- finiteness is essential for the recurrence
proof of P-indivisibility, which holds on every finite reversible
representative with a non-deterministic visible step; the representation
itself is universal, though the effective finiteness result ({[}GR,
§8.4{]}) shows the deep hidden sector may be infinite without affecting
the leading-order transition data on accessible timescales. Second, a
Turing machine is generally irreversible (it can erase, overwrite, and
halt); φ is a bijection --- nothing is erased, nothing is created, there
is no halting. Third, a Turing machine computes an \emph{extrinsic
function} (input → output for an external user); (S, φ) computes no
extrinsic function --- it is a closed permutation that cycles through
states and returns. The appearance of dynamics, probability, particles,
and forces is entirely the observer's perspective --- what the
permutation looks like through the partial window V.

\textbf{Remark (the Turing connection under effective finiteness).} The
three differences soften under the effective finiteness result. With the
deep hidden sector potentially infinite, the first difference is between
the formal definition (S finite) and the physical requirement for
accessible-timescale content (only
\(\mathcal{C}_V \times \mathcal{C}_B\) need be finite). The second
difference is a specialization, not an opposition: reversible Turing
machines (Bennett, Fredkin-Toffoli) are a well-studied subclass. The
third difference --- extrinsic vs.~intrinsic --- is the one that does
physical work. The mapping is then structural: V is the head, H is the
tape, φ is a reversible transition function, and C1--C4 characterize the
architecture. The framework extends Turing's question: instead of asking
what a machine can compute for an external observer, it asks what
computation looks like to a component of the machine --- and proves the
answer admits the quantum representation universally --- the
memory-bearing sector, where the distinctly quantum phenomenology lives,
characterized by C1/C3/C4.

\textbf{The arrow of time.} The substratum has no arrow of time --- φ
and φ⁻¹ are equally valid. Observers satisfying C1--C4 are nonetheless
structurally restricted to the non-equilibrium phase of
\((S, \varphi)\): the theorem in {[}Main, §4.6{]} shows, under its
effective-mixing hypothesis and for the coarse-grained description
generated by iterating the one-step effective channel, that
\(\tau_B(V)\) is bounded by the local mixing time on equilibrium
microstates, so C2 fails there wherever that mixing time is short
compared with \(\tau_S\) --- and hence C4, by the erasure lemma of
{[}Main §2.3{]}, which is the condition the characterization registers.
This replaces the low-entropy ``past hypothesis'' with a structural
observer-selection rule and excludes Boltzmann brains structurally
rather than anthropically. Entropy increase is then emergent as a
coarse-grained description of the non-equilibrium dynamics observers
necessarily see.

\textbf{Access, inference and computation.} A proper instantaneous
readout and a retained sequence are different observational carriers.
Gödel's and Turing's undecidability results concern their own
formal/computational hypotheses; they do not imply that a known finite
reversible model has an unidentifiable hidden state. Finite record
fibres and finite future-trace equivalence can be computed by
enumeration or partition refinement. The finite-law representation and
hidden-predictive-memory theorems do not require every hidden
distinction to remain unknown.

\textbf{Hidden-state inference.} Fix known dynamics, an admissible
initial support and the records the observer may retain. A hidden
quantity is identified precisely when it is constant on the states
compatible with those records {[}Proposition 24.5{]}. For example,
\(\varphi(v,h)=(h,v)\) with readout \(v\) has \((X_0,X_1)=(v_0,h_0)\):
two retained readings identify the initial pair and the current hidden
bit \(H_1=X_0\). With a uniform initial law it nevertheless has
\(X_2=X_0\) and positive history readback. C4 therefore need not leave
residual hidden-state uncertainty after the full record. By contrast, an
exactly decoupled spectator can remain in a non-singleton record fibre.
Cardinality comparisons across different realizations are another
question: Theorem 24(iii) preserves its specified visible channels under
normalized independent decoupled enlargement; alphabet comparisons
retain their stated hypotheses. Neither comparison proves universal
state opacity.

\textbf{Mathematics and physics.} The trace-out performs a
Jordan-Chevalley projection ({[}SM, Appendix A{]}): it extracts the
semisimple part of the dynamics and erases the nilpotent monodromy.
Physics is the semisimple shadow of mathematics --- the diagonalizable
spectral data, projected by the trace-out and organized by the gauge
group's representation structure. Theorem 23 below gives conditional
local lattice/gauge outputs on the selected cubic branch. It does not
establish an inverse reconstruction of that branch or a complete
observational gauge classification; Bell-inclusive existence and
uniqueness impose further obligations.

\begin{center}\rule{0.5\linewidth}{0.5pt}\end{center}

\hypertarget{the-reconstruction-theorem}{%
\subsection{3. The Reconstruction
Theorem}\label{the-reconstruction-theorem}}

The forward construction associates a specified reversible model and
observer interface with accessible laws {[}Main; SM §§3.1, 4--5{]}. The
converse question is which model features observations identify. The
three stages below organize possible temporal realizations and the
additional premises of a selected spatial branch. Theorem 23 aggregates
conditional outputs of that branch. It does not prove a unique inverse
within A1--A6, and H-Bell separately conditions a common Bell-compatible
completion.

\hypertarget{inputs-to-the-reconstruction}{%
\subsubsection{3.1 Inputs to the
reconstruction}\label{inputs-to-the-reconstruction}}

The reconstruction takes two kinds of inputs: empirical observations and
structural assumptions. Being explicit about both is essential to the
uniqueness claim.

\textbf{Empirical inputs} (facts about the observed universe):

(E1) \textbf{Unitary quantum mechanics.} The observed physics is
described, to the precision tested, by quantum mechanics: states as
vectors in a complex Hilbert space, unitary time evolution, self-adjoint
observables, Born-rule outcome statistics.

(E2) \textbf{Bell violations.} The observed correlations violate Bell
inequalities, ruling out local hidden-variable theories with
factorizable distributions.

(E3) \textbf{Finite boundary entropy.} The entropy of any bounded region
is finite and scales as the area of its boundary. This is a theoretical
input --- the holographic bounds {[}1, 2{]}, which are not themselves
observed --- applied to the cosmological horizon, whose finite area
makes the bound apply.

(E4) \textbf{Spatial isotropy.} Observed physics is rotationally
invariant; no spatial direction is preferred.

(E5) \textbf{Propagating gravity.} Gravitational degrees of freedom
propagate (gravitational waves are observed). This requires
\(d \geq 3\), since the Weyl tensor vanishes for \(d \leq 2\). (Used in
Stage 2(a) below; one of the three independent forward filters that
select \(d = 3\).)

(E6) \textbf{Stable matter.} Atoms, planets, and bound states exist on
observed timescales. This requires \(d \leq 3\), since the Coulomb
potential in \(d\) dimensions gives unstable atoms for \(d \geq 4\)
(Ehrenfest {[}10{]}) and gravitational orbits are unstable for
\(d \geq 4\) (Tangherlini {[}11{]}). Without stable matter, no embedded
observers exist to define the partition. (Used in Stage 2(a); second
forward filter for \(d = 3\).)

(E7) \textbf{Boundary-entropy concordance.} The cosmological parameter
inferred from observation under the standard cosmological model
satisfies \(\rho_s / \rho_{\text{crit}} \approx 1\) (spatial flatness
near critical density); it is an inferred quantity, not a directly
measured one. In \(d\) spatial dimensions, the general-\(d\) form of
this ratio is \(\rho_s / \rho_{\text{crit}} = 2/(d-1)\), which equals
unity only for \(d = 3\). (Used in Stage 2(a); third forward filter for
\(d = 3\). See {[}SM §3.2{]} and {[}GR §7.2{]} for the derivation.)

\textbf{Structural assumptions} (restrictions on the class of candidate
substrates \((S, \varphi)\)):

(A1) \textbf{Finiteness.} The configuration space \(S\) is finite.
(Two-part status. E3, with the holographic bound read as a Hilbert-space
dimension cutoff \(\dim \mathcal{H} \leq e^{A/4}\), bounds what
observation reaches: the visible sector and its boundary layer,
\(\mathcal{C}_V \times \mathcal{C}_B\) --- on which the emergent
observables depend at leading order on accessible timescales, by the
boundary-only dependence lemma {[}GR §3.2{]}. For a deep hidden sector
that no dynamics couples to the boundary layer, the extension beyond the
observer's access is observationally redundant for the specified visible
channel, rather than an established full all-partitions gauge
identification: by the deep-sector enlargement freedom (Theorem
24(iii)), such a sector may be finite of any size, or infinite, without
changing any transition probability, and A1 selects the minimal ---
finite --- representative of the equivalence class. A deep sector
coupled to the boundary layer with slow dynamics changes the transition
probabilities at \(\mathcal{O}(t/\tau_B)\); its size is not resolved by
the leading-order transition law on accessible timescales, which is the
scope of the transfer in the corollary following Theorem 24. In the
decoupled case the substrate's total cardinality beyond the boundary
layer is not an observable question, in the same sense that the uniform
vacuum offset is not.)

(A2) \textbf{Determinism.} \(\varphi: S \to S\) is a bijection
(deterministic, reversible dynamics). (Two-part status, per the dilemma
argument of {[}Main §2.2{]}: a non-injective descent either leaves a
statistical trace --- excluded where the observed statistics are of the
doubly stochastic class that the uniform-prior marginal of a bijection
produces --- or leaves none, in which case it is removable, every map on
a finite set restricting to a bijection on its eventual image, where all
recurrent registered history lives. Bijectivity for the dynamics
governing registered history is thus partly structural and partly
anchored by one empirical input, the observed statistics in that scope;
Main's honest-accounting ledger records the same division.) The anchor
is empirically live: its current frontier is nanoparticle matter-wave
interference at macroscopicity \(\mu = 15.5\) --- sodium clusters above
\(1.7 \times 10^5\) amu held in 133-nm path superpositions {[}4{]} ---
an order of magnitude beyond the prior benchmark, and the direct
one-sided test against the intrinsic-collapse alternatives that would
break it.

(A3) \textbf{Bounded coupling degree.} Each site is coupled to a bounded
number of neighbors in \(\varphi\)'s action. This restricts the graph
class used for locality and dimensional structure. The graph comparisons
of Theorem 24(iv) retain their stated hypotheses and structural
conclusions; they do not establish equality of complete channel families
for every bounded-degree graph in a quasi-isometry class.

(A4) \textbf{Center independence.} The dynamics \(\varphi\) does not
depend on a choice of ``center'' site, in two senses used below:
\(\varphi\) commutes with lattice translations up to gauge, and the
update at a site contains no explicit copy of that site's present value
({[}SM §4.1{]}). (The second sense is what the wave-equation derivation
of Stage 2 uses. Its observational anchor is the homogeneity of physical
law --- no experiment has found a preferred location in the laws --- the
cosmological principle's complement to E4's isotropy; whether it is
independently warranted by the observer architecture remains the
question flagged in §1.)

(A5) \textbf{Linearity.} The wave equation for \(\varphi\) is linear.
(Equivalent to amplitude-scale gauge invariance --- that the field-value
scale is unphysical --- via the linearity-equivalence lemma of {[}SM
§4.1{]}; hence necessary \emph{given} that gauge principle and ``partly
derived'' in exactly that sense. Whether the \(q\)-size gauge freedom of
{[}SM §2.7{]} entails amplitude-scale gauge is an identified open step;
pending it, nonlinear alternatives are excluded precisely to the extent
amplitude-scale gauge is assumed, and would otherwise require a separate
derivation.)

(A6) \textbf{Background independence.} The dynamics is covariant under
spatially-varying internal-index transformations: a site-dependent
transformation \(G(\mathbf{n})\) of the internal index is a symmetry of
the dynamics provided the coupling data carried on the links are
transformed with it,
\(M(\mathbf{n}, \hat{e}_j) \to G(\mathbf{n})\,M(\mathbf{n}, \hat{e}_j)\,G(\mathbf{n}+\hat{e}_j)^{-1}\)
({[}SM §3.1{]}). The content of the assumption is this covariant
interface --- the coupling is link-valued data on which the
transformations act --- and once it is in place the local transformation
law imposes no further condition on the rule. A6 is not the stronger
fixed-background condition that the same transformations leave the
dynamics invariant with the coupling held fixed and preserved pointwise:
for an invertible constant coupling that condition forces the
transformation to agree across every coupled pair of sites, and at the
symmetric point of the coupling matrix it fails for every transformation
that differs across a coupled pair. The global commutant symmetry ---
the transformations constant across the lattice that leave the coupling
itself unchanged, the specialization of A6 in which nothing needs
transporting --- is what the local gauge reading of {[}SM §3.1{]}
proceeds from. The state-dependent coupling graph of {[}SM §3.1{]},
under which the graph itself evolves with the state, is a separate
principle that shares the name and is not A6.

Theorem 23's branch-specific conclusions use E1--E7 and A1--A6 together
with additional model-selection premises named in §3.3; E1--E7 and
A1--A6 alone are not shown to identify a unique substratum. The set
A1--A6 is used in the conditional local construction but is not proved
sufficient to select a unique realization, nor proved minimal: A4
(center independence) and A6 (background independence) overlap in
physical content (A6 promotes the symmetry that A4 constrains), and A5
(linearity) is equivalent to amplitude-scale gauge invariance via the
linearity-equivalence lemma of {[}SM §4.1{]}, hence partly derived ---
necessary \emph{given} amplitude-scale gauge, with whether {[}SM
§2.7{]}'s \(q\)-size gauge entails it left as an identified open step. A
tighter axiomatization may be possible but is not pursued here; the
stated local consequences also depend on the explicit branch choices and
completion premises, not only on the independence of A1--A6. The
empirical inputs E5, E6, E7 (propagating gravity, stable matter,
boundary-entropy concordance) provide explicit forward filters for
\(d = 3\) --- see Stage 2(a) below.

\textbf{Critical dependencies.} Theorem 23's proof relies on the
following prior theorems, whose individual correctness is assumed:

\begin{itemize}
\tightlist
\item
  {[}Main §2.3{]} history readback for the \textbf{memory-bearing
  temporal} sector used in Stage 1 (a modeling premise about the
  temporal data, not an inference from Bell violation)
\item
  {[}Main §3.3{]} Bell ceiling and the adopted measurement-independent,
  ontically parameter-dependent Bell branch (a separate
  spatial-composite premise; Bell violation is not used as a witness of
  temporal P-indivisibility)
\item
  {[}Main §3.2{]} Forward Stinespring dilation (bijection → CPTP
  channel) and reverse-direction lemma (one-step doubly stochastic
  matrix → uniform-prior bijection; exact for rational entries, to
  arbitrary precision otherwise)
\item
  {[}Main §3.4{]} Finite-horizon process dilation (stochastic process →
  bijection, per finite horizon)
\item
  {[}Main §3.4{]} Characterization theorem (accessible non-Markovianity
  --- the framework's finite-horizon visible-law memory notion, not
  generic quantum non-Markovianity --- ⟺ C1/C3/C4 per finite accessible
  horizon; C2's slow-bath necessity conditional on the hidden-sector
  mixing hypothesis, ETH its physical motivation, its persistence form
  holding without it)
\item
  {[}SM §3.2{]} Coupling-graph dimension → d = 3
\item
  {[}SM §4.1{]} Center independence + isotropy + linearity → wave
  equation
\item
  {[}SM Theorems 5--15{]} Gauge group, generations, hypercharges
  derivation chain
\item
  {[}SM Theorems 17--21{]} T-invariance → \(\theta\in\{0,\pi\}\) and
  reciprocal transition counts; \(\bar\theta\) open (H-top, H-det)
\item
  {[}GR §3{]} Gap equation from thermal self-consistency →
  \(\hbar = c^3 \epsilon^2 / (4G\sigma)\)
\end{itemize}

Theorem 23's branch-relative conclusions inherit the stated limitations
of these dependencies and of the additional assumptions in §3.3.

\hypertarget{stage-1-observed-qm-deterministic-embedding-with-c1c4}{%
\subsubsection{3.2 Stage 1: Observed QM → deterministic embedding with
C1--C4}\label{stage-1-observed-qm-deterministic-embedding-with-c1c4}}

\textbf{Inputs:} E1 (unitary QM), E2 (Bell violations), E3 (finite
boundary entropy), A1 (finiteness), A2 (determinism), and M1, in two
clauses: \textbf{M1-T}, the observed temporal quantum data used in Stage
1 lie in the accessible history-readback sector of {[}Main
§2.3--§3.4{]}; and \textbf{M1-B}, Bell-violating composites take {[}Main
§3.3{]}'s measurement-independent, ontically parameter-dependent branch
with operational no-signaling. M1-T and M1-B are logically distinct;
Bell violation is not taken to imply temporal P-indivisibility.

\textbf{Output:} There exists a triple \((S, \varphi, V)\) with \(S\)
finite, \(\varphi\) a bijection on \(S\), and \(V \subset S\) a
distinguished subset such that: - The observed quantum dynamics is the
reduced description obtained by tracing out \(S \setminus V\) from the
deterministic evolution under \(\varphi\). - The triple satisfies C1
(non-trivial coupling), C2 (record persistence --- the slow bath one
realization), C3 (sufficient hidden memory capacity), and C4 (history
readback) for the M1-T temporal sector: {[}Main §3.4{]}'s universal
hidden-memory theorem --- accessible non-Markovianity ⟺ C1/C3/C4
realization per accessible horizon, with \(M_t > 0\) forcing unavoidable
hidden predictive memory --- makes readback-relevant hidden memory
unavoidable in every faithful deterministic realization. M1-B separately
requires Bell-relevant cross-wing coupling/parameter dependence and does
not supply C4.

\textbf{Conditional structure.} C1, C3, and C4 for the memory-bearing
temporal sector are derived from E1--E3, A1--A2 together with
\textbf{M1-T} --- C4 via the hidden-memory theorem applied to the
stipulated accessible temporal memory; conditional on M1-T, they hold
without further hypothesis. \textbf{M1-B} is separate: it maps E2 to the
ontically parameter-dependent Bell composite and supplies an independent
cross-wing-coupling requirement, not C4. C2 takes the two forms of
{[}Main §1.3{]}. (C2-structural), retention of the record of visible
history to its readback, holds by reversibility in every bijective
realization. C2's necessity holds in its persistence form without
further hypothesis: in the coupling-step model of {[}Main §3.4{]},
observed conditional memory \(I^* > 0\) places some realized hidden
conditional at the coupling time at least
\(\varepsilon_{\min}(I^*, n_V)\) from its reference prior (the
physical-memory theorem). Its slow-bath form is conditional on the
quantitative mixing hypothesis of {[}Main §3.4{]} --- uniform
total-variation closeness of the hidden sector's post-relaxation
conditionals to a single bath distribution --- under which the same
bound makes \(\varepsilon_{\min}(I^*, n_V)\) a lower bound on the
relaxation residue at the coupling time, hence on the bath's conditional
mixing time, and a departure \(\delta\) from the homogeneous Markov
benchmark at step order \(k\) forces \(\varepsilon \geq \delta/(k-1)\)
(the bound \(\|T^{(k)} - T^k\|_{\mathrm{TV}} \leq (k-1)\varepsilon\),
constant optimal at \(k = 2\)); without it, accessible conditional
memory does not bound the bath's mixing time. ETH is the physical case
for the hypothesis, as discussed in {[}Main §3.4{]}'s ``Remark (Status
of ETH).'' ETH is a well-supported conjecture for generic chaotic
many-body Hamiltonians and is the standard assumption for the
cosmological-horizon complement; relaxations or failures of ETH would
require a separate argument. The framework's gauge-group derivation
through Stage 2 inherits this conditional.

\textbf{Bell--gravity compatibility hypothesis (H-Bell).} The Stage-2
reference nearest-neighbor wave-equation graph has a pointwise causal
cone advancing one edge per update; whether that cone or the lift's
long-wavelength speed \(1/\sqrt3\) is the physical \(c\) is not derived
{[}Main §3.3{]}. Under measurement independence, setting interventions
that each read and write only their own wing and \(\pm1\) readouts
outside the remote setting's graph cone, {[}Main §3.3{]} therefore makes
that \textbf{reference graph} Bell-local and caps CHSH at \(2\). The
framework nevertheless already makes the graph state-dependent: the
discrete Einstein equation of {[}SM §3.1{]} is formulated on \(G(x)\),
and its mod-\(q\) spatial-Markov proof explicitly assumes uniform
initial conditions. M1-B may therefore be implemented by
\textbf{preparation-indexed adjacency}: the vacuum/uniform reference
class retains the geometric bounded-range graph used by the local SM/GR
derivation, while an entangling preparation supplies bounded-degree
Bell-relevant long edges (or an equivalent ontically nonlocal response
map) furnishing parameter dependence without operational signaling.

H-Bell is the remaining compatibility statement for that branch: the
preparation-indexed graphs must (i) supply the required ontic parameter
dependence, (ii) preserve operational no-signaling, and (iii) satisfy a
metric/Ollivier--Ricci stability condition strong enough that the
continuum-curvature step of {[}SM §3.1{]} remains valid on the state
classes to which the Einstein reconstruction is applied. The
graph-boundary area law is not the primary unresolved issue:
\(|\partial V|\) is defined on \(G(x)\) and therefore counts
preparation-induced crossing edges. The open step is the identification
of that graph boundary/curvature with continuum geometric area/Ricci
curvature. {[}SM §3.1{]}'s sparse local-contamination estimate,
\(L_Nm_N/N\to0\), is a sufficient condition for typical mesoscopic
neighborhoods to remain reference-like, but it does not prove global
curvature convergence, control defect endpoints, or establish the
integrated-curvature limit. H-Bell is therefore an explicit
curvature-stability/composite-existence hypothesis, not a consequence of
A1--A6 or of the nearest-neighbor wave equation.

Two points limit the exposure this conditional creates. First, the
relevant question is not whether ETH holds universally --- it is known
to fail for integrable and many-body-localized (MBL) systems --- but
whether it holds for \emph{the} hidden sector implied by the framework's
other commitments. The hidden sector here is the high-capacity
complement (\(N_H \gg N_V\) --- the comfortable realization of C3's
capacity floor) of a bounded-coupling bijection on a statistically
isotropic graph (A3, E\_isotropy). Integrability and MBL are both
non-generic in a specific sense relevant here: integrability requires an
extensive set of local conserved quantities, which a statistically
isotropic bounded-degree coupling graph without fine-tuned commuting
structure does not possess; MBL requires strong quenched disorder
sustaining localization against the thermalizing tendency of generic
interactions, and is moreover known to be unstable in dimension
\(d \geq 2\) (the avalanche instability). Since the reconstruction
independently selects \(d = 3\) (§3.3, {[}SM §3.2{]}), the hidden sector
sits in the regime where MBL does not survive and generic interacting
dynamics thermalize --- so ETH is the expected case, not an additional
fine-tuning. Second, the conditional attaches only to the slow-bath
(timescale) form of C2's necessity direction; C2's persistence form and
the sufficiency direction --- a C1--C4 realization yields a visible law
with conditional memory (forced by C4), and that finite law universally
admits the quantum representation --- are unconditional, so the
framework's positive predictive content does not rest on ETH. The mixing
hypothesis, for which ETH is the physical case, is used only to convert
observed memory into a lower bound on the bath's conditional mixing
time; the reconstruction's uses of that form hold within the class of
hidden sectors satisfying the mixing hypothesis. ETH, which the argument
above makes the expected case at \(d = 3\), is the physical case for
that hypothesis and does not by itself supply it: a hidden sector
satisfying ETH for few-body observables while retaining
conditional-distribution memory violates it ({[}Main §3.4{]}, Remark
(role of ETH)).

\textbf{Derivation.} The observed quantum dynamics provides, for each
finite experimental context, a stochastic law on a finite visible
configuration space. By {[}Main{]}'s exact finite-horizon equivalence,
this law has a finite reversible deterministic realization and an exact
fixed-basis unitary/Born representation. The representation step is
universal and therefore does not itself reconstruct the nontrivial OI
sector. The additional modeling premise M1 separates two empirical
structures. \textbf{M1-T} places the temporal data used for Stage 1 in
the accessible history-readback sector; {[}Main §2.3{]} then proves that
on a fixed finite representative a non-permutation rooted map, which C4
readback supplies, forces indivisibility somewhere in the recurrence
cycle. \textbf{M1-B} places the Bell data on {[}Main §3.3{]}'s
measurement-independent, ontically parameter-dependent branch; Bell
violation therefore certifies Bell-relevant cross-wing coupling but is
not used to infer temporal P-indivisibility. By {[}Main{]}'s
finite-horizon equivalence (\(S \Rightarrow D\), the process-dilation
and response-table constructions of {[}Main §3.4{]}), the temporal law
admits a realization as the visible marginal of a deterministic
bijection \(\varphi\) on an enlarged finite space
\(\mathcal{C}_V \times \mathcal{C}_H\), with a hidden prior independent
of the visible initial state carried as part of the realization datum;
for a single one-step doubly stochastic matrix the reverse-direction
lemma ({[}Main §3.2{]}) gives the uniform-prior form, exactly when its
entries are rational and to arbitrary precision otherwise. This delivers
the triple \((S, \varphi, V)\) with
\(S = \mathcal{C}_V \times \mathcal{C}_H\) and \(V = \mathcal{C}_V\).
Non-permutation \(T(t)\) in the M1-T temporal sector gives the
corresponding C1 witness; independently, M1-B's ontic parameter
dependence requires Bell-relevant cross-wing coupling. The necessity
analysis of {[}Main §3.4{]} establishes, for every faithful realization
of the resulting dynamics, C3 (the data-processing bound
\(\log_2|\mathcal{C}_H| \geq I^*\), with no mixing hypothesis) and
(C2-structural) (retention of the record of visible history to its
readback, supplied by reversibility; record persistence is required for
the non-Markovian returns that produce quantum interference); in the
coupling-step model of {[}Main §3.4{]}, observed conditional memory
keeps some realized hidden conditional at least
\(\varepsilon_{\min}(I^*, n_V)\) from the reference prior at the
coupling time (C2's persistence form, with no mixing hypothesis), and
under the hidden-sector mixing hypothesis ({[}Main §3.4{]}; ETH its
physical motivation) this lower-bounds the bath's conditional mixing
time (the slow-bath form).

\textbf{Stage-1 dilation freedom.} For the finite channel/dilation data
used in temporal realization, hidden-basis freedom is defined at the
level of an individual channel. Exact hidden-factor relabelling and
independently decoupled finite enlargement preserve the full
time-resolved channel under the hypotheses of Theorem 24. Neither fact
classifies all dilations of a common channel family. This is not
uniqueness of the full triple \((S,\varphi,V)\) across all coherent
composites. In particular, it does not select a preparation-indexed Bell
adjacency rule, and the operational lift still requires its separately
stated finite-predictive-dimension, local-tomography and purification
hypotheses.

\textbf{Status:} Finite-horizon realization and its stated memory
conditions; no unique dilation, selected spatial branch or
Bell-inclusive completion follows from this stage alone.

\hypertarget{stage-2-embedding-isotropy-specific-lattice-gauge-structure}{%
\subsubsection{3.3 Stage 2: Embedding + isotropy → specific lattice +
gauge
structure}\label{stage-2-embedding-isotropy-specific-lattice-gauge-structure}}

\textbf{Inputs:} Stage 1 output + E4 (spatial isotropy) + E5
(propagating gravity) + E6 (stable matter) + E7 (boundary-entropy
concordance) + A3 (bounded coupling) + A4 (center independence) + A5
(linearity) + A6 (background independence) + max-lattice-speed condition
(for the wave-equation coupling constant; see step (b) below).
Identification of this local lattice output with the Bell-violating
Stage-1 completion additionally assumes H-Bell.

\textbf{Output:} The local propagating substratum sector is a \(d = 3\)
cubic lattice bijection with: - \(K = 2d = 6\) internal components per
site - Coupling matrix eigenvalue multiplicities \((3, 2, 1)\) - Gauge
group \(\text{SU}(3) \times \text{SU}(2) \times \text{U}(1)\) - Three
generations of chiral fermions in the SM representations - One Higgs
doublet \((\mathbf{1}, \mathbf{2}, +1/2)\) - Anomaly-free hypercharges
\((Y_Q, Y_u, Y_d, Y_L, Y_e) = (1/6, 2/3, -1/3, -1/2, -1)\)

\textbf{Derivation.}

\begin{enumerate}
\def\labelenumi{(\alph{enumi})}
\item
  \emph{Dimension.} The coupling graph inherited from Stage 1 has
  polynomial growth exponent \(d\). Three independent forward filters
  from the empirical inputs converge on \(d = 3\): propagating gravity
  (E5) requires \(d \geq 3\) (the Weyl tensor vanishes for
  \(d \leq 2\)); stable matter (E6) requires \(d \leq 3\) (Coulomb
  instability of atoms for \(d \geq 4\), plus Tangherlini's
  gravitational orbit instability); boundary-entropy concordance (E7)
  gives \(\rho_s / \rho_{\text{crit}} = 2/(d-1) = 1\) at \(d = 3\)
  uniquely. The intersection of the first two filters already pins
  \(d = 3\); the third is a sharpening consistency check that the
  framework passes exactly. The integer-polynomial growth assumption
  (rather than fractal or exponential growth) is itself derived from
  self-consistency: exponential growth fails Lorentz-invariant
  dispersion (Jacobson's GR derivation requires it); fractal growth
  fails the cubic-lattice symmetry needed for the cubic-group
  decomposition in step (d) below. By Gromov's theorem and statistical
  isotropy (E4), the surviving graph class is quasi-isometric to
  \(\mathbb{Z}^d\) for integer \(d\); the three filters then select
  \(d = 3\). ({[}SM §3.2{]} presents this argument in full; the inputs
  E5, E6, E7 surfaced here make explicit the empirical content the
  filters require.)
\item
  \emph{Wave equation.} Center independence (A4, in its functional
  sense: no self-term), isotropy (E4), and linearity (A5) uniquely
  select, within the nearest-neighbor class (an assumption beyond A3's
  bounded degree; see the scope remark at {[}SM §4.1{]}), the wave
  equation up to a multiplicative coupling constant \(\alpha\) ({[}SM
  §4.1, Theorem{]}): the unique second-order linear dynamics on a
  lattice that is translation-invariant, isotropic, reversible and free
  of a self-term has the form
  \(f = \alpha(x_1 + x_2 + \cdots + x_{2d}) \bmod q\) with coupling
  coefficient \(\alpha\). The coefficient is neither the propagation
  speed nor free to be set to \(1\): causal support travels at most one
  lattice edge per update whatever its value, and on the observer-level
  normalized branch it equals \(1/d\) ({[}SM §4.1, Remark{]}).
\item
  \emph{Internal components.} H-link gives \(K=2dm\), with the
  additional single-copy clause \(m=1\) giving \(K=6\); coupling-degree
  minimization does not establish that clause {[}SM §4.7{]}. The
  selected staggered construction has a singlet/triplet candidate-sector
  decomposition. The identification of the triplet as three physical
  fermion generations retains H-spin' and the separate matter
  hypotheses; it does not follow from the link count alone.
\item
  \emph{Gauge group.} The six geometric links decompose as
  \(T_1\oplus E\oplus A_1\). Under H-link single-copy + H-cust, a
  generic equivariant condensate has unitary stabilizer
  \(G_0=U(3)\times U(2)\times U(1)\) {[}SM Theorems 5 and 7{]}. The
  pointwise commutant of the cubic action itself is \(U(1)^3\), distinct
  from this stabilizer. A precise correlated two-torus quotient gives
  \(G_0/K\cong S(U(3)\times U(2))\) {[}SM Proposition 14a{]}, and the
  selected SM charged family has faithful image
  \((SU(3)\times SU(2)\times U(1))/\mathbb Z_6\) {[}SM Proposition
  15a{]}. Amplitude-scale invariance alone does not select \(K\) or that
  matter action; the full native endomorphism carrier has only the
  common scalar circle in its kernel. The physical action and remaining
  phase selection are additional premises. A6 supplies the separate
  covariance of site-dependent transformations with transformed link
  data, not the global quotient.
\item
  \emph{Matter content.} The selected staggered construction gives a
  singlet and three degenerate triplet sectors; their physical
  spin/generation interpretation retains H-spin', H-χ' and
  light-spectrum selection {[}SM Theorems 8--13{]}. The link carrier
  admits cross-charged matter representations without selecting the
  physical family. For the family fixed in {[}SM Theorem 15{]}, anomaly
  cancellation gives the complete charge classification, including the
  neutral-\(Q\) exception. Nonzero \(Y_Q\) recovers the SM ratios up to
  normalization and quark-singlet interchange; physical charge
  normalization remains separate.
\item
  \emph{Discrete symmetries.} T-invariance of the bijection \(\varphi\)
  gives \(\theta\in\{0,\pi\}\) and reciprocal visible transition counts
  at every scale (Theorems 17--21); it does not fix \(\bar\theta\),
  which remains open under H-top and H-det {[}SM §5{]}. Reciprocity
  leaves the Jarlskog invariant free; only a mechanism making \(Y_u\)
  and \(Y_d\) simultaneously real would force \(J=0\), which is the
  falsification control on H-det proposals {[}SM, §5.3--5.5{]}.
\end{enumerate}

\textbf{Scope of the Stage-2 selection.} Steps (a)--(f) do not supply a
chain of converse implications from the listed outputs to a unique
substratum. The dimension filters use E5--E7; step (b) works within a
nearest-neighbor, second-order, linear wave-rule class stronger than
bounded coupling alone; step (c) uses H-link and the selected
multiplicity \(m=1\); and the subsequent representation and matter
conclusions use the cubic/octrahedral representation, amplitude-scale
gauge and the qualifications in {[}SM{]}. These give conditional
structural outputs on the selected branch. They neither exclude every
other substratum satisfying E1--E7 and A1--A6 nor establish uniqueness
modulo a fully characterized \(\mathcal{G}_{\text{sub}}\).

\textbf{Bell compatibility.} Step (b)'s nearest-neighbor
\textbf{reference} cone cannot itself realize M1-B, while retaining
measurement independence, in a Bell protocol whose wings are separated
by more than the graph cone over its duration. Stage 2 therefore
establishes the local lattice/gauge sector on the geometric reference
class. The state-dependent \(G(x)\) machinery leaves a native route for
M1-B through preparation-indexed adjacency, but the Ollivier--Ricci
continuum theorem presently cited by {[}SM{]} does not cover those
long-edge perturbations. Nor does it cover the unperturbed reference
graph, and there the situation is worse than open: the degree-\(6\)
cubic hop metric is exactly \(\ell_1\), with a scale-independent
\(\sqrt2\) diagonal stretch, so the hop-metric curvature route fails for
the reference family and no \(D=3\) extension of the cited theorem would
rescue it. The repair is to re-found curvature on the
propagation/Laplacian geometry, which the substratum's isotropic leading
dispersion already supplies {[}SM §3.2{]}. Promoting the local sector
and M1-B to one common \((S,\varphi)\) is therefore conditional on
H-Bell's curvature-stability requirement.

\textbf{Status:} Conditional results for the selected local lattice
branch; uniqueness over all candidate substrata and a common
Bell-inclusive completion remain open.

\hypertarget{stage-3-dynamics-emergent-fundamental-constants}{%
\subsubsection{3.4 Stage 3: Dynamics → emergent fundamental
constants}\label{stage-3-dynamics-emergent-fundamental-constants}}

\textbf{Inputs:} Stage 2 output.

\textbf{Output:} The emergent action-scale relation
\(\hbar = c^3 \epsilon^2 / (4G\sigma)\), equivalently
\(\epsilon = 2\sqrt{\sigma}\,l_p\), where \(l_p\) is the Planck length
and \(\sigma\) the entropy per boundary mode ({[}GR §3.1{]}); it fixes
only \(\epsilon/\sqrt{\sigma}\), not the value of \(\hbar\) ({[}GR
§4{]}).

\textbf{Derivation.} The gap equation from thermal self-consistency of
the boundary layer ({[}GR §3{]}) gives the stated relation. The
derivation uses: existence of thermal equilibrium at the boundary layer
(a consequence of the C1--C4 dynamics from Stage 1), the boundary-only
dependence lemma ({[}GR §3.2{]}), and the infrared detailed-balance
condition H-slope with the horizon and frame conditions ({[}GR §2, §3.2,
§8.5{]}), which does not follow from \(\varphi\) being a bijection
({[}GR §3.2{]}) and is an assumption beyond Stages 1--2.

\textbf{Uniqueness at Stage 3.} The gap equation admits a unique
solution under the stated inputs ({[}GR §3.2{]}).

\textbf{Status:} Conditional on H-slope with the horizon and frame
conditions ({[}GR §3.2{]}).

\hypertarget{the-composite-theorem}{%
\subsubsection{3.5 The composite theorem}\label{the-composite-theorem}}

\textbf{Lemma 23.0} (Shared outputs on the selected local branch).
\emph{Suppose two finite reversible models both implement the same
selected local branch: spatial dimension \(d=3\), a nearest-neighbor
cubic wave-rule carrier, H-link with \(m=1\), the octahedral
representation and the amplitude-scale-gauge and matter-sector
hypotheses explicitly invoked in {[}SM{]}. Then they have the same
values of the branch's discrete representation invariants \(K=6\) and
multiplicities \((3,2,1)\), together with those gauge and matter outputs
that follow from the cited {[}SM{]} results under their respective
hypotheses. If both additionally satisfy H-slope and the horizon/frame
assumptions of {[}GR §3.2{]}, they obey the same action-scale relation
\(\hbar=c^3\epsilon^2/(4G\sigma)\); this does not fix a common numerical
\(\hbar\) without fixing \(\epsilon\) and \(\sigma\). Equality of these
outputs does not imply the two substrata are related by any proposed
gauge generator.}

\emph{Proof.} Under the common selected branch and H-link with \(m=1\),
\(K=2dm=6\); the stipulated octahedral representation has the fixed
decomposition \((3,2,1)\). The remaining stated gauge and matter outputs
are inherited, only under their individual hypotheses, from the
corresponding {[}SM{]} results. The {[}GR{]} calibration under H-slope
supplies the relation among \(\hbar\), \(\epsilon\) and \(\sigma\). None
of these deductions identifies the full microscopic dynamics or
reconstructs a unitary from the single-time visible channel; Theorem 24
establishes only its stated fixed-interface invariances and Proposition
24.2 has stronger block-word hypotheses. \(\square\)

\textbf{Theorem 23} (Conditional local-branch reconstruction).
\emph{Assume the finite reversible framework with the stated E1--E7 and
A1--A6 inputs, together with a selected \(d=3\) nearest-neighbor cubic
wave-rule realization, H-link at \(m=1\), the requisite octahedral
representation, amplitude-scale gauge and the separate matter-sector
premises in {[}SM{]}. On that selected branch the internal carrier count
is \(K=6\) and its octahedral multiplicities are \((3,2,1)\); the stated
gauge-group, matter, chirality and hypercharge outputs follow under
their respective {[}SM{]} hypotheses, and the relation
\(\hbar=c^3\epsilon^2/(4G\sigma)\) follows only under H-slope and the
horizon/frame assumptions in {[}GR{]}. The branch's reference local
causal dynamics cannot by itself supply the Bell-violating M1-B
composites under measurement independence; a common completion is
conditional on H-Bell. The hypotheses do not establish uniqueness of the
local branch among all candidate substrata, completeness of a substratum
gauge group, or uniqueness of a Bell-inclusive completion.}

\emph{Proof.} The selected wave-rule and H-link premises give the
internal-count and octahedral outputs by the local constructions of §3.3
and {[}SM{]}. The further gauge, matter and hypercharge statements use
exactly the additional premises and derivation steps named in {[}SM{]};
they are not promoted to unconditional consequences of E1--E7 and
A1--A6. Under the separate {[}GR{]} horizon conditions the calibration
relation follows. Stage 1's finite dilation represents temporal laws;
this alone supplies neither the selected spatial branch nor a
Bell-compatible common realization. The reference nearest-neighbor
causal graph is Bell-local over separated wings under the stated
measurement-independence and readout conditions; the proposed
preparation-indexed completion is subject to H-Bell and is not
constructed here. Lemma 23.0 compares only the stipulated branch's
outputs, and neither it nor Proposition 24.2 classifies all substrata
modulo gauge. \(\square\)

\hypertarget{remarks}{%
\subsubsection{3.6 Remarks}\label{remarks}}

\textbf{Remark} (Evidential weight of the construction). The selected
cubic branch yields representation multiplicities \((3,2,1)\) under
H-link with \(m=1\) and the stipulated octahedral action. Gauge and
matter interpretations inherit their separate premises, including
H-cust, H-spin' and H-χ' where invoked; the hypercharge result retains
its family-pattern assumptions. {[}SM §7{]} records classified
quantitative comparisons within that branch. Agreement with Standard
Model data tests the applicable chains, but does not establish that
A1--A6 and the empirical inputs select this branch uniquely. The
physical branch selector and the observational equivalence
classification require independent arguments.

\textbf{Remark} (Conditional content of the local construction). Theorem
23 states consequences of explicit premises. Once the \(d=3\)
nearest-neighbor cubic wave-rule branch, H-link with \(m=1\) and its
octahedral representation are fixed, \(K=6\) and the multiplicities
\((3,2,1)\) follow. Further gauge, generation, chirality and hypercharge
conclusions use the respective {[}SM{]} physical hypotheses; the
action-scale relation additionally uses H-slope and the horizon/frame
conditions of {[}GR{]}. These deductions constrain a chosen
construction. They do not show that all models compatible with E1--E7
and A1--A6 share these structures or belong to one gauge orbit. Equality
of the listed outputs also does not identify the full microscopic
dynamics. The inverse selection problem remains open within A1--A6
without further premises.

\textbf{Remark} (Derived outputs, empirical selection and model
premises). The representation multiplicities follow from the selected
link representation. The physical interpretation of the resulting blocks
requires the separate carrier and matter hypotheses of {[}SM{]}. The
empirical arguments for \(d=3\) and the cubic crystal system are
conditional selection arguments in their stated classes; they do not
classify every reversible model allowed by A1--A6. The nearest-neighbor
wave rule, site factorization and single-copy H-link clause are
additional premises. The derivation-status discussion in {[}SM §8.3{]}
and the observable classifications in {[}SM §7{]} distinguish these
roles. None supplies a theorem that observations recover a unique local
branch or a complete microscopic gauge class.

The structural assumptions A1--A6 exclude several broad classes of
theory from the outset:

\begin{itemize}
\tightlist
\item
  Intrinsically continuum theories (violate A1)
\item
  Intrinsically stochastic theories (violate A2)
\item
  Theories with unbounded coupling degree (violate A3)
\item
  Theories violating center independence, linearity, or background
  independence (violate A4--A6)
\end{itemize}

Within the selected cubic branch, the specified local outputs follow
under the hypotheses listed in Theorem 23. Uniqueness of all
realizations modulo a complete observational gauge group is not
established, and neither is Bell-inclusive substratum uniqueness. A
separate argument would be required to rule out alternatives outside
this class. The framework's defense of A1--A6 (in {[}SM §2{]} and
elsewhere) is that they are natural given the empirical inputs E1--E7,
but this defense is itself an argument, not a theorem.

\textbf{Remark} (Hypothesis dependencies). Each structural assumption
can in principle be weakened or replaced:

\begin{itemize}
\tightlist
\item
  A1 (finiteness) is supported by E3; alternative would be an
  effective-finiteness argument with continuum UV completion, which the
  framework does not develop.
\item
  A2 (determinism) is the key ontological commitment; stochastic
  alternatives would produce different reconstructions.
\item
  A3 (bounded coupling) is conventional for physical systems; long-range
  couplings would give different dimensional structure.
\item
  A4 (center independence) rules out preferred-frame theories.
\item
  A5 (linearity) is the strongest restriction; it is equivalent to
  amplitude-scale gauge invariance (linearity-equivalence lemma, {[}SM
  §4.1{]}), so nonlinear wave equations on finite lattices are the
  separate class one obtains exactly when that gauge principle is
  dropped.
\item
  A6 (background independence), as covariance under local internal-index
  transformations with the link coupling transformed alongside, is the
  standard gauge-theoretic requirement; the stronger fixed-background
  invariance, with the coupling held fixed, is a separate condition and
  is not what A6 asserts.
\end{itemize}

Weakening any \(A_i\) does not merely enlarge the equivalence class; it
potentially produces \emph{different} reconstructions that may not agree
with observation. The programme proposes that E1--E7 + A1--A6, together
with the selected cubic branch and further conditions, form an
explanatory model; a unique selection of that model and the exclusion of
all alternatives have not been established.

\textbf{Remark} (Continuum extension). The lattice-level predictions are
the framework's primary claims --- the lattice is the fundamental
description, not an approximation to a continuum theory. The structural
results (gauge group, representations, generation count) are
algebraic/topological properties of the lattice theory and are exact.
Quantitative observables (scattering amplitudes, mass ratios) are
compared to experiment through continuum perturbation theory; the
lattice-continuum discrepancy at any experimentally accessible energy
\(E\) is suppressed by
\((E\epsilon/\hbar c)^2 = 4\sigma(E/M_{\text{Pl}})^2\), which at
\(\sigma = 1\) is \(\sim 10^{-32}\) at LHC energies. The rigorous proof
that lattice Yang-Mills defines a continuum limit with a mass gap (a
Clay Millennium Prize problem) is an open problem in mathematics, but it
is not required by the OI framework: the lattice theory is complete, and
its predictions are lattice-exact.

The epistemic structure is as follows. A reasonable objection is that
the framework treats the lattice as fundamental (so the mass-gap problem
can be set aside) while also using continuum perturbation theory to
extract the numbers that match data (which would make the continuum
limit load-bearing). The position is consistent. The framework's claim
is that the lattice theory is the complete physical content and its
predictions are exact \emph{at the lattice level}. The comparison to
experiment is then a comparison between two objects: the lattice
theory's prediction extracted via continuum PT (a controlled
approximation with error bounded by
\((E\epsilon/\hbar c)^2 = 4\sigma(E/M_{\text{Pl}})^2 \sim 10^{-32}\) at
\(\sigma = 1\)), and the continuum-QFT fit to the experimental data (the
standard way the data is reported). The empirical agreement is therefore
agreement between the lattice theory's continuum approximation and the
continuum-QFT data fit, with the approximation error in the first object
\(\sim 30\) orders of magnitude below the experimental precision of the
second, so the lattice-vs-continuum gap is not the source of any
observed agreement or disagreement at accessible energies. The mass-gap
problem is not load-bearing for this: it concerns whether a continuum
limit \emph{exists as a rigorous mathematical object}, whereas the
framework requires only that continuum PT is an accurate
\emph{calculational approximation} to the lattice theory's predictions
at low energy --- the standard status of lattice-to-continuum
perturbative matching, which does not require the existence of the
rigorous limit. The lattice is fundamental, the continuum is a
calculational tool with a quantified error, and the experimental
comparison is between the tool's output and continuum-reported data.

\textbf{Remark} (Finite observation sets vs.~idealized empirical
inputs). The empirical inputs E1--E7 are stated as structural facts
about observed physics --- the full apparatus of QM, Bell violations
attaining Tsirelson's bound, the holographic entropy bound, exact
rotational invariance, propagating gravitational waves, stable matter on
all observed scales, and spatial-curvature concordance. Operationally,
no finite observation set fully establishes any of these; finite
measurements support them only in the sense of being consistent with the
data so far. Theorem 23 gives conditional outputs for the selected cubic
branch even when the empirical inputs are idealized as exact; neither
local-branch selection nor Bell-inclusive uniqueness follows from those
idealized observations alone. Reconstruction from finite observation
sets is a related but distinct question: with finite data, multiple
equivalence classes may all be consistent, with the data discriminating
among them only weakly. The framework's operational answer to the
finite-data question is the cumulative weight of the quantitative
predictions of {[}SM §7{]} --- twenty-two classified retrodictions of
gauge couplings, CKM/PMNS angles, mass ratios, and the Koide relation
(two parameter-free structural and five conditional structural, the
remainder layered, chained or fitted, the gauge couplings among the
fitted), the measured comparisons each at \(\lesssim 1\sigma\) --- which
collectively rule out reconstructions that would have produced different
values. Theorem 23 states the branch-specific structural consequences,
while {[}SM §7{]} records experimental comparisons for quantities
computed or fitted within that branch; those comparisons alone do not
prove uniqueness of its microscopic realization.

The implication supplied by Theorem 23 is conditional on the selected
cubic realization and its named physical premises:

\[\text{Selected branch and hypotheses}\quad\Longrightarrow\quad\text{local outputs}.\]

The arrow does not give a converse from observations to the selected
branch. A physical selector, an appropriate observational equivalence
relation and a valid classification would be needed for a stronger
inverse theorem; the restricted coherent-channel generators already fail
to be exhaustive by Pair C. H-Bell separately conditions a common
Bell-compatible completion, whose uniqueness is unproved. Interpreting a
model's mathematics as reality or as a description produces the same
measurements only when the interpretations stipulate the same accessible
data. That factoring statement does not turn a complete ontology or a
resolution of Wigner's puzzle into a consequence of Theorem 23 or
Theorem 24.

\textbf{Remark} (ER=EPR at substratum and emergent levels). The
Maldacena-Susskind ER=EPR conjecture {[}3{]} proposes that any two
entangled systems are connected by an Einstein-Rosen bridge, with the
strong form asserting that entanglement \emph{is} a non-traversable
wormhole --- the same ontological structure under two descriptions. The
framework here engages this conjecture at two distinct levels.

At the \textbf{emergent level}, the conjecture's \emph{weak} form is
reproduced: any pair of entangled subsystems in the emergent QM
description shares boundary modes via the partition trace-out, and the
holographic dictionary maps shared boundary modes to a connecting
geometric structure in the emergent gravitational description. This
recovers the ER=EPR weak form --- entangled systems are connected by a
wormhole-like geometric object in the emergent theory --- as a
structural consequence of the framework's holographic architecture
rather than as an independent postulate.

At the \textbf{substratum level}, the conjecture's \emph{strong} form
(entanglement is identical to wormhole geometry, not merely correlated
with it) is predicted \emph{false} at the representative level. Two
distinct partition representatives can produce identical emergent EPR
correlations while having different substratum representative structures
--- there is no faithful bijection from EPR pairs to substratum
geometric features. The strong identity holds at most up to gauge
equivalence on the substratum side. Whether a gauge-orbit reformulation
of strong ER=EPR --- ``EPR-equivalence-classes correspond to
ER-equivalence-classes'' --- is true in the framework remains open; the
test would be whether the substratum's gauge group
\(\mathcal{G}_{\text{sub}}\) acts transitively on the set of substratum
representatives that produce a given EPR correlation pattern.

This produces a sharp empirical signature: the framework's operational
predictions agree with weak ER=EPR (which would also be predicted by
AdS/CFT, by quantum error correction in holography, and by other
emergent-geometry approaches) but disagree with strong ER=EPR at the
representative level. Distinguishing these experimentally would require
probing substratum-scale features, which is currently beyond reach. The
disagreement is therefore a structural prediction rather than a testable
one at present, but it concretely positions the framework relative to
the broader ER=EPR program.

\begin{center}\rule{0.5\linewidth}{0.5pt}\end{center}

\hypertarget{companion-family-identification}{%
\subsubsection{3.7 Companion-family
identification}\label{companion-family-identification}}

Within the specified finite companion family \(\Phi_F(x,y)=(F(x)-y,x)\),
a known labelled finite abelian group and exact complete current-field
triple support, from a full-support phase ensemble, determine \(F\).
Under uniform phase counting the classical one-step population kernel is
instead uniform for every \(F\); this need not hold for a structured
stationary prior. The
\href{../verification/audits/manuscript/theorem-recovery-20261028/native-companion-identification.md}{companion
identification note} gives the proof, a complete-parent-patch version
and a coarse-readout control. This restricted inverse assumes its
readout and ensemble coverage; a single orbit need not supply them. It
supplies no physical branch selector, and Theorem 23's inverse selection
problem remains open.

\begin{center}\rule{0.5\linewidth}{0.5pt}\end{center}

\hypertarget{the-substratum-gauge-group}{%
\subsection{4. The Substratum Gauge
Group}\label{the-substratum-gauge-group}}

For a specified visible/hidden decomposition \(S=V\times H\) and a
specified hidden preparation \(\rho_H\), the \textbf{fixed-interface
coherent-channel equivalence} \(U\sim_{\mathrm{ch}}U'\) means equality,
for every integer \(t\geq0\) and every visible input density matrix
\(\rho\), of the channels \[
\Phi_t^U(\rho)=\operatorname{Tr}_H[U^t(\rho\otimes\rho_H)(U^\dagger)^t].
\] The weaker \textbf{population equivalence} \(U\sim_{\mathrm{pop}}U'\)
compares only
\(T_{ij}(t)=\langle i|\Phi_t^U(|j\rangle\langle j|)|i\rangle\) for every
\(i,j,t\), at the same interface and preparation. A third reading
compares observables \textbf{across all admissible observer partitions};
it requires a specified correspondence between those partitions, their
preparations and their readouts. No classification of that third
relation is proved here. A simultaneous relabelling of the dynamics
\emph{and} all prepared/read-out observables is only a passive change of
description; it cannot certify equivalence at a fixed interface. The
expression \(\mathcal G_{\mathrm{sub}}\) below denotes a
\textbf{proposed} observational gauge structure relative to a stated
observable carrier, not an already classified common gauge group for
these three readings.

\textbf{Theorem 24} (Exact interface invariances and their different
scopes). \emph{Let \(V,H\) be finite and let \(U\) be a permutation
unitary on \(\mathbb C^V\otimes\mathbb C^H\), with uniform hidden
preparation. (i) A single hidden permutation \(\sigma:H\to H\), applied
identically in every visible fibre, preserves \(\Phi_t\) for all \(t\)
under conjugation by \(\mathbb1_V\otimes P_\sigma\). More generally, a
fibre-dependent hidden relabelling \(\pi(v,h)=(v,\sigma_v(h))\)
preserves all rooted transition probabilities \(T_{ij}(t)\), but need
not preserve coherent channels. (iii) Adjoining an exactly decoupled
ancillary dynamics and an independently prepared normalized ancilla
preserves \(\Phi_t\) for all \(t\) wherever the tensor-product channel
and partial trace are defined. Clauses (ii) and (iv) below are separate
structural/conditional comparisons, not assertions of equality of all
coherent channels. Restricted hidden-only relabelling and decoupled
enlargement do \textbf{not} generate every fixed-interface
coherent-channel equivalence (Pair C below). Completeness across all
partitions remains open.}

\emph{(i) Relabellings and observation.} For a single \(\sigma\), its
permutation unitary is \(\mathbb1_V\otimes P_\sigma\), and the maximally
mixed hidden state is invariant. The outer hidden conjugations disappear
under the partial trace, so the coherent-channel identity holds for
arbitrary visible density matrices. For \(\pi(v,h)=(v,\sigma_v(h))\),
the visible component of \(\pi\) is unchanged, and each rooted fibre's
uniform hidden prior is permuted bijectively. Therefore the visible
output counts, at every iterate, are unchanged. However the controlled
hidden permutation generally acts differently on off-diagonal visible
coherences and need not leave \(\Phi_t\) invariant. An
\textbf{arbitrary} permutation of all states need preserve neither data
set when the readout is fixed. For example, take \(V=H=\{0,1\}\),
\(\varphi(v,h)=(h,v)\), and exchange \((0,0)\) with \((1,0)\) while
fixing the other two states. At \(t=1\) the rooted transition matrices
are \[
T_\varphi(1)=\begin{pmatrix}1/2&1/2\\1/2&1/2\end{pmatrix},
\qquad
T_{\pi\varphi\pi^{-1}}(1)=\begin{pmatrix}1&0\\0&1\end{pmatrix}.
\] These disagree already on populations and thus on the coherent
channels. If the readout and preparation are also transported by
\(\pi\), the agreement is a passive coordinate identity, not a new gauge
invariance. A visible-layer swap \(T\) in a specific wave-rule branch
need not be of hidden-only form and is \textbf{not} thereby a generator
in the coherent-channel group; its possible time-reversal interpretation
is a separate structural statement.

\emph{(ii) Alphabet size.} The results in {[}SM §2.7{]} that hold for
every \(q\geq2\) transfer \textbf{without} an extra alphabet-size
hypothesis. Other quantities in that section retain explicit \(q\)
dependence. Equality of transition probabilities across different
alphabet sizes is a \textbf{separate hypothesis}, and the unconditional
\(q\)-insensitive structural conclusions do not establish that
hypothesis or equal coherent-channel families.

\emph{(iii) Decoupled ancillary enlargement.} For a separate system
\(D\), a normalized ancillary density operator \(\rho_D\), and
unitary/bijective independent dynamics \(\psi\), the enlarged evolution
is \(U^t\otimes\psi^t\). Its reduced visible channel obeys \[
\operatorname{Tr}_{H\otimes D}[(U^t\otimes\psi^t)
(\rho_V\otimes\rho_H\otimes\rho_D)
(U^\dagger{}^t\otimes\psi^\dagger{}^t)]
=\Phi_t^U(\rho_V)\operatorname{Tr}(\rho_D)=\Phi_t^U(\rho_V).
\] This is exact for finite \(D\) and also for infinite-dimensional
\(D\) with well-defined unitary dynamics, normalized trace-class
\(\rho_D\) and the indicated partial trace. It does not transfer
finite-state recurrence to an infinite enlargement. A \textbf{coupled}
deep sector, however slowly coupled, is not covered by this equality;
the accessible-time approximation has the separate error term of {[}GR
§3.2{]}.

\emph{(iv) Structural graph comparisons.} Theorems 1--6 and {[}SM
§2.2{]} retain their \textbf{specified} structural conclusions within
the graph classes and assumptions to which those derivations apply:
polynomial-growth dimension, mode-count and statistical-isotropy
conditions can support corresponding emergent-QM, dimension and
gravitational \emph{structural} content. This is stronger than simply
noting common graph labels, but it is \textbf{not} an equality theorem
for complete rooted channels. The Standard Model cubic decomposition
additionally uses the octahedral point-group representation; generic
quasi-isometry does not preserve it. Neither these structural results
nor a graph isomorphism alone proves coherent-channel equivalence at a
fixed observer interface.

\emph{Proof of the exact invariances.} The coherent statement in (i)
follows from conjugation by the single hidden unitary, invariance of the
uniform hidden density, and the partial-trace identity. For the
population statement in (i), each \(\sigma_v\) bijectively permutes the
preparation fibre \(H\) and fixes the final visible coordinate; applying
\(\pi\varphi^t\pi^{-1}\) proves equality of every rooted count.
Statement (iii) follows from tensor-product evolution, unitary
preservation of the independent ancillary trace, and partial trace. The
structural claims (ii) and (iv) retain only the indicated source
hypotheses and conclusions. \(\square\)

\textbf{Completeness depends on the observable carrier.} For
\textbf{fixed-interface coherent channels} \(\sim_{\mathrm{ch}}\), the
proposed restricted hidden-only transformations and decoupled
enlargements do \textbf{not} generate all equivalences: the six-state
Pair C dynamics have equal \(\Phi_t\) for every \(t\) but are separated
by the balanced block-word trace invariant tested by \texttt{ST3\_pairC}
and \texttt{ST4\_pairC}. For finite ancillary carriers the normalized
trace invariant survives hidden-only conjugation and independent
decoupled enlargement; its behavior under an allowed visible-coordinate
relabelling is also included in the tested Pair C obstruction. Thus
merely chaining these moves cannot connect the pair. This is a
\textbf{restricted non-exhaustiveness result}, not a statement about
arbitrary additional moves. For \textbf{population-only} equivalence
\(\sim_{\mathrm{pop}}\), the larger fibre-dependent relabellings
preserve \(T_{ij}(t)\) and the Pair C obstruction does not establish
non-exhaustiveness under that relation. For \textbf{all-partitions}
equivalence, the relevant observational information and permitted
correspondences have not been settled, so completeness remains
\textbf{open}. None of these three verdicts can be substituted for
another.

A written operator certificate extends the restricted Pair C obstruction
to arbitrary nonempty ancillary Hilbert spaces, without requiring an
infinite-dimensional normalized matrix trace. For the visible blocks
define
\[W_0=B_{00}B_{10}B_{01}B_{00}^{\dagger}B_{01}^{\dagger}B_{10}^{\dagger},\]
and let \(W_1\) interchange the two visible labels. Exact multiplication
gives \((W_0,W_1)=(\operatorname{diag}(1,0,0),0)\) for Pair C and
\((0,0)\) for the primed Pair C. Hidden-unitary conjugation conjugates
each word; adjoining independent unitary dynamics replaces each balanced
word by \(W_i\otimes I_D\). Zero versus nonzero is therefore unchanged
for any nonempty finite or infinite \(D\), and a visible-label swap only
exchanges the two statuses. The unordered status pair separates the
examples through every finite chain of these moves and their inverses.
This operator argument is written mathematics; the cited Lean ancillary
certificate retains its finite-carrier scope. It supplies no
all-partitions classification or invariance under additional moves.

\textbf{Corollary (effective finiteness and gauge-class transfer).}
Three statements make A1's two-part status precise. \emph{(i) The
bound.} E3, read as a Hilbert-space dimension cutoff, bounds the
observationally relevant space:
\(|\mathcal{C}_V \times \mathcal{C}_B| \leq e^{S_\partial}\) with
\(S_\partial = A/4\) in Planck units --- for the de Sitter horizon,
\(e^{10^{122}}\). This dimension-cutoff reading is the interpretive
premise of A1, and this is the only place it enters. \emph{(ii)
Boundary-only dependence.} For accessible times \(t \ll \tau_B\),
\(T_{ij}(t) = T_{ij}^{(B)}(t) + \mathcal{O}(t/\tau_B)\) ({[}GR §3.2{]}),
so at leading order every emergent observable on accessible timescales
is a functional of \(\mathcal{C}_V \times \mathcal{C}_B\) data alone.
\emph{(iii) Transfer.} Under a decoupled enlargement of the kind in
Theorem 24(iii), with a normalized independent ancillary state and
well-defined partial trace, every \(T_{ij}(t)\) is unchanged at every
\(t\), whether the ancilla is finite or infinite-dimensional. Statements
expressible solely through those probabilities transfer; finite-state
recurrence and other statements not fixed by that law do not transfer to
infinite systems. An extension coupled to the boundary layer with
\(\tau_B^D \gg \tau_S\) alters \(\mathcal{C}_D\) and its dynamics while
leaving \(T^{(B)}\) unchanged, so the two systems share only the
leading-order law: a statement expressible through
\(\{T^{(B)}_{ij}(t) : t \ll \tau_B\}\) transfers by (ii), while a
statement about the transition probabilities themselves transfers only
if it is insensitive to \(\mathcal{O}(t/\tau_B)\) corrections. A
transition probability equal to zero on one system and to
\(\varepsilon = \mathcal{O}(t/\tau_B) > 0\) on the other is the
elementary statement that does not transfer. P-indivisibility at the
recurrence scale --- established on finite representatives by recurrence
(\(\varphi^N = \mathrm{id}\), {[}Main §2.3{]}) --- is not an
accessible-timescale statement: across coupled extensions it holds on
the finite members, and on separately specified infinite-deep-sector
models, where finite recurrence is unavailable, the readback lemma
({[}Main §2.3{]}) gives the accessible-window property it states,
conditional memory, a different and weaker conclusion. A1's choice of
the finite representative is therefore a convenience of proof for
content carried by the leading-order accessible-timescale law;
recurrence-scale content is stated for finite representatives.
\(\square\)

\textbf{Remark (amplitude-scale gauge).} The operational principle used
here declares an overall field normalization redundant when preparations
and readouts are compared under the corresponding alphabet-coordinate
change. Affine changes \(\phi\mapsto\lambda\phi+c\), with \(\lambda\) a
unit, are available coordinate relabellings; observational invariance
for a specified carrier is an additional requirement, not a conclusion
about every fixed readout. The principle is separate from alphabet-size
comparisons and from Theorem 24's hidden-only channel invariance. Its
role in the conditional linearity argument of {[}SM §4.1{]} does not
imply that independent block phases can all be removed while retaining
hypercharge. The exact correlated-quotient construction and its
additional physical-action premise are stated in {[}SM Proposition
14a{]}.

\emph{Status (open, referee-grade).} Amplitude-scale gauge is adjoined
as an explicit operational principle: it does not follow from Theorem
24(i), which here allows only hidden-only channel relabellings (and, for
population data, fibre-dependent hidden relabellings); clause (ii)
concerns alphabet-\emph{size} comparisons, not arbitrary field-value
changes. Results invoking amplitude-scale gauge therefore remain
conditional on that separate principle. Stated operationally it is
\emph{not} viciously circular: the input (absolute amplitude is
unobservable) is strictly weaker than the output (additive dynamics),
and a nonlinear \(F\) on the same alphabet is a coherent object the
principle \emph{excludes} rather than presupposes. The one route that
would make linearity \emph{unconditional} --- deriving the
unobservability of absolute amplitude from the observer architecture
itself --- is by contrast circular: what the embedded observer can
resolve through the trace-out is dynamics-dependent (a nonlinear \(F\)
induces amplitude-dependent emergent dispersion, which \emph{is}
observable), so ``absolute amplitude is unobservable'' already
presupposes linearity. Amplitude-scale gauge therefore has the status of
an operational input on a par with center independence and isotropy, not
a theorem; whether it is independently warranted by the observer
architecture is the question flagged for external review.

\emph{Formalization status of the linearity step (resolved as a
sharpened axiom, not a closure).} A dedicated audit of this step settles
its logical status, and the outcome is worth stating precisely because
it is the prototype for how the framework's ``open'' foundational steps
resolve. The linearity-equivalence lemma --- \emph{all first
finite-differences of the update rule over \(\mathbb{Z}/q\mathbb{Z}\)
are base-point-independent iff the rule is affine as a function} --- has
been verified exhaustively for \(q \le 6\) in one neighbour-coordinate
and \(q \le 3\) in two, and on \(2.5\times10^5\) random rules at larger
\((q,m)\), with zero counterexamples; together with the telescoping
proof (constant single-coordinate differences force
\(F(x)=F(0)+\sum_i c_i x_i\), well-defined on the torus since
\(q c_i \equiv 0\)) this is a complete elementary result, not a
genericity heuristic. The question of whether the \(q\)-\emph{size}
gauge freedom of {[}SM §2.7{]} \emph{entails} amplitude-scale gauge
resolves in the negative, and sharply: the weak reading of §2.7
(size-independence of predictions) does not reach rescaling at fixed
\(q\); the strong reading (arbitrary relabellings of
\(\mathbb{Z}/q\mathbb{Z}\) are gauge) \emph{overshoots} --- a nonlinear
permutation conjugating a linear rule preserves cycle structure (hence
emergent transition probabilities up to relabelling) while being
nonlinear, so it would render a nonlinear rule gauge-equivalent to a
linear one and thereby \emph{destroy} linearity-necessity rather than
establish it (verified explicitly on \(\mathbb{Z}/5\mathbb{Z}\)). The
principled object that delivers exactly what is needed without
overshooting is the \textbf{additive-automorphism} group
\(x \mapsto \lambda x + c\) (\(\lambda\) a unit) --- a proper subgroup
of all permutations that preserves the alphabet's ``+'', hence the
dispersion. The upshot, in the framework's three-tier classification
({[}SM §8.3{]}): linearity is \textbf{not} promoted to a forced (Tier-1)
consequence; it is a Tier-3 \emph{definitional stipulation}, but a
sharpened one --- the bare assumption ``the dynamics is linear'' is
replaced by the more primitive and better-motivated assumption ``the
alphabet's additive automorphisms are gauge,'' from which linearity
follows as a certified theorem. This is the characteristic resolution
shape for the framework's foundational gaps: pressing for a proof does
not close the gap to a derivation but relocates the irreducible
assumption to a more primitive and more defensible place, and names it
exactly.

\textbf{Classification scope of the transformations.} Theorem 24
establishes hidden-only coherent-channel invariance and finite or
trace-class decoupled enlargement, and identifies a broader
population-only hidden-fibre relabelling. Pair C refutes exhaustiveness
of the \emph{restricted coherent-channel} moves; the all-partitions
classification remains open. Equality of all time-resolved visible
channels does not determine the bijection up to hidden unitary, even
with trivial hidden commutants.

Specht's trace criterion and Wiegmann's simultaneous-unitary extension
characterize unitary similarity through traces of words in matrices and
their adjoints {[}7--9{]}. Proposition 24.2 is a fixed-carrier
specialization with the additional spanning premise used by its cited
Lean witnesses. The classical finite-dimensional criterion itself does
not require that premise.

\textbf{Proposition 24.2} (Fixed-carrier block-word rigidity). \emph{Let
\(V,H\) be finite index sets and \(U,U'\) be matrices on
\(\mathbb C^V\otimes\mathbb C^H\). Write \(B_{vw}(U)\) for their visible
blocks on \(\mathbb C^H\), and let \(\mathcal W(U)\) be the linear span
of all words in the blocks and their adjoints, including the identity.
Suppose \(\mathcal W(U)=M_{|H|}(\mathbb C)\) and the unnormalized trace
of every corresponding block-word agrees for \(U\) and \(U'\). Then a
hidden unitary \(W\) satisfies
\(U'=(\mathbb 1_V\otimes W)U(\mathbb 1_V\otimes W)^\dagger\).
Conversely, hidden-unitary conjugation preserves every block-word trace,
and the hypotheses imply \(\mathcal W(U')=M_{|H|}(\mathbb C)\).} No
unitarity assumption on \(U,U'\) is needed, but the two systems must
have the same hidden index set.

\emph{Proof.} Equality of all word traces yields a trace-preserving
unital \(*\)-isomorphism of the generated word spans. When the first
span is the full matrix algebra, the matrix-unit construction implements
it by a unitary. Necessity follows by trace invariance. The Lean
witnesses for the two directions are WT2gen\_oneSided and
trace\_wordEval\_hiddenConjugate, combined in WT3\_spanning\_iff;
blockSpanning\_of\_wordTracesEqual establishes spanning for the target
(OIBridge/WordTraceSufficiency.lean and
OIBridge/SemigroupTransfer.lean). \(\square\)

\textbf{Limits of the channel-family data.} In the Lean-checked
ST3\_pairC counterexample, two permutation unitaries on six states
(\(|V|=2\), \(|H|=3\)) induce the same visible channel \(\Phi_t\) at
every time, with nontrivial visible statistics and trivial hidden
commutants. Their traces are \(2\) and \(0\), so no unitary conjugates
them. ST4\_pairC also excludes the tested same-size
decoupled-enlargement and hidden-relabelling alternative. They have
unequal block-word traces (pairC\_wordTraces), so they do not meet
Proposition 24.2's stronger premises. These counterexamples refute the
channel-family-to-hidden-conjugation implication and, under the
restricted fixed-interface moves, generator exhaustiveness. They do not
classify unrestricted population or all-partitions equivalence.

\textbf{Open scope questions.} Applying Proposition 24.2 to
reconstruction requires the selected observables to determine all
block-word traces and the relevant dynamics to have a full block span. A
hidden unitary also need not be a finite-state relabelling of form (i).
The fixed-interface coherent-channel generator proposal is refuted by
Pair C under the restricted moves above; classification of population
equivalence and equivalence over all admissible partitions remains open.
Whether additional coherent or interleaved multi-time data are
physically available is an independent question.

\hypertarget{conditional-reconstruction-from-complete-word-data}{%
\subsubsection{4.1 Conditional reconstruction from complete word
data}\label{conditional-reconstruction-from-complete-word-data}}

Multi-time records can recover more than separate one-time channels, but
their observational carrier must be specified. The following classical
minimal-realization result for weighted automata and rational series
originates in Schützenberger, Carlyle--Paz and Fliess {[}10--12{]};
minimal-realization uniqueness up to similarity is stated in {[}5,
Chapter 2, Theorem 2.4{]}, with the common sequential-system and Hankel
framework developed in {[}6{]}. This result and its native deterministic
specialization recover predictive structure. They complement Proposition
24.2, whose input is coherent block-word traces, rather than imply its
hypotheses.

\textbf{Proposition 24.3} (Minimal linear word realizations). \emph{Let
a finite alphabet \(\mathcal A\) have two finite-dimensional
realizations \((E,a,b,(M_c))\) and \((E',a',b',(M'_c))\) over a field,
with chronological convention
\(M_{c_1\cdots c_k}=M_{c_k}\cdots M_{c_1}\) and \(M_\epsilon=I\). Assume
each is reachable, \(\operatorname{span}\{M_wa\}=E\), and observable,
\(\bigcap_u\ker(bM_u)=\{0\}\), and similarly for the primed model.
Equality \(bM_wa=b'M'_wa'\) for every word implies a unique linear
isomorphism \(S\) with \(Sa=a'\), \(SM_c=M'_cS\), and \(b'S=b\). For
arbitrary competing dimensions \(n,m\ge1\), equality through word length
\(n+m-1\) already implies equality for all words. A nonzero series of
known minimal dimension \(d\) has a scalar Hankel reconstruction using
prefix and suffix lengths at most \(d-1\) and word statistics through
\(2d-1\).}

\emph{Proof.} Define \[
S\!\left(\sum_i t_iM_{w_i}a\right)=\sum_i t_iM'_{w_i}a'.
\] If the source sum is zero, every target continuation evaluates the
proposed image to
\(\sum_i t_i b'M'_uM'_{w_i}a'=\sum_i t_i bM_uM_{w_i}a=0\); target
observability makes the image zero. Thus the map is well-defined.
Reversing the construction gives its inverse by reachability. The empty
word gives \(Sa=a'\), appending a letter gives \(SM_c=M'_cS\), and
equality of evaluations gives \(b'S=b\). Any such map agrees on the
reachable generators, proving uniqueness.

For the length bound, the reachable spaces
\(R_k=\operatorname{span}\{M_wa:|w|\le k\}\) satisfy
\(R_{k+1}=R_k+\sum_cM_cR_k\). Equality at two consecutive stages makes
the space invariant and the sequence stable; otherwise its dimension
increases. In dimension \(N\) it therefore stabilizes by \(N-1\) for
nonzero \(a\), and the zero-vector case is immediate. Apply this to the
\((n+m)\)-dimensional discrepancy realization with operators
\(M_c\oplus M'_c\), vector \((a,a')\), and evaluation \((b,-b')\).

In a minimal \(d\)-dimensional model choose reachable-basis prefixes
\(w_j\) and observable-basis suffixes \(u_i\), each of length at most
\(d-1\) by the same filtration argument. Put \(R=[M_{w_j}a]_j\) and
\(O=[bM_{u_i}]_i\). Then \[
H_{ij}=s(w_ju_i)=(OR)_{ij},\qquad
(H_c)_{ij}=s(w_jcu_i)=(OM_cR)_{ij}.
\] Both \(O\) and \(R\) are invertible. The reconstructed coordinates
are \[
\widehat M_c=H^{-1}H_c,\qquad
\widehat a=H^{-1}(s(u_i))_i,\qquad
\widehat b=(s(w_j))_j.
\] These equal \(R^{-1}M_cR\), \(R^{-1}a\), and \(bR\). The largest
required length is \((d-1)+1+(d-1)=2d-1\), and the invertible minor
proves that full Hankel rank is \(d\). \(\square\)

The alphabet must be stationary and freely composable, with every
required prefix, suffix and extension accessible; an outcome-labelled
instrument uses an unnormalized linear update. A restricted protocol
grammar or changing apparatus requires its own state augmentation or
admissible-language theorem. The dimension/rank bound is an independent
hypothesis: a finite observed Hankel rank alone does not certify it.
Similarity is only a linear conclusion. Even strictly positive
normalized minimal stochastic models can have a unique nonmonomial
similarity with an inverse having negative entries. An isomorphism
mapping the entire normalized classical simplex onto itself would
instead permute its extreme points, but generic linear reachability and
observability do not establish that stronger condition.

\textbf{Proposition 24.4} (Passive finite-permutation predictive
quotient). \emph{Let \(P\) be a permutation of a nonempty finite set
\(\Omega\), with deterministic passive readout \(r:\Omega\to\mathcal A\)
and initial law \(\mu\). Identify \(\omega,\eta\) exactly when
\(r(P^t\omega)=r(P^t\eta)\) for every \(t\ge0\), and let \(Q\) be the
resulting future-trace quotient. It inherits a permutation \(\bar P\),
readout \(\bar r\), and prior \(\bar\mu\). For each model being
compared, assume its quotient prior \(\bar\mu\) has full support, and
let \(B\) be a known upper bound on both quotient-state counts. For
\(B\ge2\), equality of all passive word probabilities through length
\(B\), counting the time-zero readout, determines
\((Q,\bar P,\bar r,\bar\mu)\) up to a unique readout-preserving
conjugacy transporting the prior. For \(B=1\), one-letter data suffice.
No invariance or uniformity of \(\mu\) is required. The quotient's
linear Hankel rank equals \(|Q|\); common-length retained-word vectors
and block effects generate its full nonnegative state cone and all
coordinate indicators, respectively. Minimal-realization similarity
between two such quotients is a permutation.}

\emph{Proof.} Future-trace equivalence is forward invariant. Since \(P\)
has a finite order, every negative power is also a nonnegative power, so
the equivalence is backward invariant too and \(\bar P\) is a
permutation. On an \(N\)-state deterministic labelled system let \(C_k\)
be the partition by its first \(k\) readouts, with \(C_0\) indiscrete.
If \(C_k=C_{k+1}\), equality of labels and successor classes makes this
equivalence forward invariant and no later refinement occurs. Otherwise
the number of classes increases. Thus \(C_{N-1}\) is the full future
partition for \(N\ge2\).

Put \(\ell=B-1\ge1\). Distinct quotient states are distinguished by
their length-\(\ell\) future blocks. Full support makes the common
positive-probability length-\(\ell\) words exactly the state labels
\(W\). Each \(w\in W\) has exactly one symbol \(x\) with \(s(wx)>0\);
its successor label is the last \(\ell\) symbols of \(wx\). The first
symbol of \(w\) gives its readout and \(s(w)\) gives its initial mass.
Hence length-\(B\) data determine the labelled quotient permutation and
prior. Any readout-preserving intertwiner preserves the entire future
trace, so this conjugacy is unique. For \(B=1\) the sole state has
identity dynamics and its one-letter law determines its label.

For the cone and rank statements choose one common separating length
\(\ell\ge\max(1,|Q|-1)\). In quotient coordinates set
\(M_x=\bar P D_x\), where \(D_x\) filters readout \(x\), \(a=\bar\mu\),
and \(b=\mathbf1^T\). These reproduce passive cylinder probabilities
because the last advance cancels under \(\mathbf1^T\bar P=\mathbf1^T\).
If \(w_q\) is the length-\(\ell\) block of state \(q\), then \[
bM_{w_q}=e_q^T,\qquad
M_{w_q}\bar\mu=\bar\mu(q)e_{\bar P^\ell q}.
\] Full support gives a positive coefficient, and the common power
\(\bar P^\ell\) bijectively covers all end states. Thus the
retained-word rays and block effects cover the coordinate rays and
indicators. The model is reachable and observable, so its Hankel rank is
\(|Q|\). Matching blocks and their probabilities makes the similarity of
Proposition 24.3 carry the corresponding coordinate rays with unit
scale, hence be a permutation. \(\square\)

\textbf{Access and remaining gauge gap.} This is a predictive-quotient
result for passive, nondisturbing records from one fixed finite
reversible carrier. It refines {[}Main §3.4{]}'s existing finite-horizon
canonical tail quotient; it does not infer an all-time carrier from
arbitrary finite-horizon laws. Full support on the microscopic carrier
is sufficient, but one rooted visible preparation
\(\delta_v\otimes\mu_H\) is generally not such a prior. A consistent
positive mixture of every root would supply it only if that preparation
family is actually accessible with the same dynamics and hidden prior.
Algebraic saturation on \(Q\) does not assert externally selectable
microscopic states or separate states in the same future-trace fibre.
Without a known state bound, full-support uniform laws on binary de
Bruijn cycles of orders \(k\) and \(k+1\) agree through length \(k\) and
differ later. Invisible fibres can also have different microscopic cycle
structures. Neither proposition therefore supplies hidden-factor
conjugacy, coherent block-word traces, a common quantum instrument
theory, or completeness over all admissible partitions.

The self-contained derivation and bounded exact rational controls are
retained in
verification/audits/manuscript/multitime-predictive-recovery-20261028/.
The native checks cover 768 four-state cases and two one-state cases,
including nonstationary full-support priors, constant readouts,
common-length cone/Hankel checks, and countercontrols for missing
hypotheses. These are executable evidence supplementing the written
proofs; neither new proposition is claimed Lean-proved.

\textbf{Proposition 24.5} (Passive record fibres and state
identification). \emph{Let \(P\) be a known permutation of a nonempty
finite carrier \(\Omega\), let \(r:\Omega\to\mathcal A\) be a
deterministic readout, and let \(\mu\) be a known initial law with
support \(A\). Suppose the observer may retain the exact passive
readings \(X_j=r(P^jS_0)\). For a positive-probability word
\(w=(w_0,\ldots,w_{k-1})\) of length \(k\ge1\), define}
\[F_w=\{s\in A:r(P^js)=w_j\text{ for }0\le j<k\}.\] \emph{(i) A quantity
\(q(S_0)\) is identified by that record exactly when \(q\) is constant
on \(F_w\). In particular the initial state is identified exactly when
\(|F_w|=1\). The current state \(P^{k-1}S_0\) has the same singleton
criterion. When \(\Omega=\mathcal C_V\times\mathcal C_H\) and
\(r=\pi_V\), the current hidden state is identified exactly under that
same criterion.} \emph{(ii) Put \(N=|\Omega|\) and \(L=\max(1,N-1)\).
Passive words of length \(L\) already induce the full future-trace
equivalence on \(\Omega\), hence on \(A\) by restriction. Every
positive-probability finite record remains compatible with at least two
initial states if and only if every nonempty intersection of a full
future-trace class with \(A\) has at least two elements.}

\emph{Proof.} Conditioning gives
\[\Pr(S_0=s\mid w)=\frac{\mu(s)\mathbf1_{F_w}(s)}{\sum_{u\in F_w}\mu(u)}.\]
Every point of \(F_w\) has positive mass, so a quantity is known exactly
when it is constant there. Invertibility maps \(F_w\) bijectively to the
current compatible-state set. For a product carrier with current visible
value recorded, distinct current compatible states have distinct hidden
coordinates. This proves (i).

For (ii), partition \(\Omega\) by its first \(j\) readings, starting
with the indiscrete partition at \(j=0\). Each next reading refines it.
If two consecutive partitions agree, the classes have identical labels
and successors in the same class; no later refinement occurs. Otherwise
the class count strictly increases. Starting with one class, there can
be at most \(N-1\) strict increases, so length \(N-1\) suffices for
\(N\ge2\); one reading suffices when \(N=1\). Restriction to \(A\) gives
the assertion without assuming that \(A\) is invariant. Each nonempty
finite-record fibre is a union of restricted full future-trace classes.
Thus classes of size at least two make every such fibre non-singleton;
conversely a singleton full class yields a singleton length-\(L\) fibre.
\(\square\)

\textbf{Memory and opacity are different.} Under the uniform law in the
swap model, \(H_1=X_0\), \(X_2=X_0\) and \(I(X_0;X_2\mid X_1)=1\), but
\(H(H_1\mid X_0,X_1)=0\). More generally, cyclically shifting \(m+1\)
independent uniform bits and reading the first gives a hidden store of
\(2^m\) states, all identified by \(m+1\) retained readings. At time
\(m\), \(I(X_{<m};X_{>m}\mid X_m)=m\) bits, saturating the hidden
capacity, while \(H(H_m\mid X_{\le m})=0\). Indeed, conditional on the
current bit, the preceding \(m\) readings specify all other bits, and
the next \(m\) readings repeat them. C4's information is predictive
memory relative to the present readout, not a lower bound on residual
uncertainty after the full history.

\textbf{Access and resource scope.} The proposition assumes the
dynamics, preparation support and retained records are available. It
supplies no physical record-storage or intervention mechanism. A stored
record with only \(R\) possible values cannot distinguish all \(|A|>R\)
admissible initial states: some fibre has at least
\(\lceil |A|/R\rceil\) members, although other fibres may be singletons.
This pigeonhole limit is weaker than claiming ambiguity for every
history. A decoupled spectator gives genuine non-singleton passive
classes; fixing that spectator by preparation, reading it, or allowing a
separating intervention can remove the ambiguity. Proposition 24.4
instead reconstructs an unknown predictive quotient from exact word laws
under its own support and state-bound hypotheses. Neither result
identifies the physical OI model, coherent instruments or the full gauge
relation. The bounded exact controls are retained in
verification/audits/manuscript/hidden-state-identifiability-20261028/.

\textbf{Other structural cases.} The \(q\)-insensitive conclusions of
{[}SM §2.7{]} are unconditional within their stated domain, while
equality of transition probabilities across \(q\) remains an additional
hypothesis. Theorems 1--6 retain their conditional structural graph
conclusions without implying equality of full channels; Standard Model
representation content requires octahedral structure. In the selected
wave-rule branch a layer swap \(T\) satisfying
\(\varphi^{-1}=T\circ\varphi\circ T^{-1}\) supplies a conditional
time-reversal relation, \textbf{not} a hidden-only instance of generator
(i). A global phase of the emergent unitary is not a separate
finite-bijection operation.

\textbf{A proposed gauge hierarchy, with an open substratum step.} The
three descriptive levels below carry distinct, stated equivalence
notions. Their nesting as groups and a complete projection of substratum
equivalences to emergent symmetries require additional hypotheses;
neither follows from Theorem 24 or Proposition 24.2 alone.

\emph{Level 3 (substratum):} At a fixed visible/hidden interface,
Theorem 24(i),(iii) exhibits exact coherent-channel invariances under a
uniform hidden-factor relabelling and any well-defined decoupled
normalized ancillary enlargement, and population invariance under
fibre-dependent hidden relabellings. Its alphabet and graph clauses
concern narrower structural comparisons, and a complete all-partitions
equivalence relation has not been characterized. The notation
\(\mathcal{G}_{\text{sub}}\) denotes the proposed full substratum
observational gauge structure only once its observables, relabellings
and partition correspondence have been specified; it is not established
here as generated by (i)--(iv).

\emph{Level 2 (emergent QFT):} The selected cubic branch supplies an
exact \((3,2,1)\) link representation and a conditional
\(U(3)\times U(2)\times U(1)\) condensate stabilizer. The SM Lie algebra
and faithful charged-family image require the specified correlated
quotient and matter action {[}SM §4.4--§4.9{]}. A complete map from
substratum observational equivalence to that physical gauge action
remains open.

\emph{Level 1 (emergent Hamiltonian):} The diagonal-unitary
transformation \(H\mapsto DHD^\dagger\) preserves the transition
probabilities in the specified measurement basis. Energy shifts and the
antiunitary reflection also preserve that data; completeness of the
stated alternatives is conditional on H-BG and the gap/overlap
hypotheses {[}GR §3.3{]}.

These are compatible descriptive levels, not a proved inclusion chain
\(\mathrm{G1}\subseteq\mathrm{G2}\subseteq\mathrm{G3}\). Establishing
such inclusions requires definitions and maps that preserve the same
observational data at each step, rather than only graph invariants or
fixed-interface channels. The substratum equivalence, unlike an emergent
Lagrangian gauge symmetry, concerns representations of a finite
bijection and specified observer access; no substratum Lagrangian is
postulated.

\begin{center}\rule{0.5\linewidth}{0.5pt}\end{center}

\hypertarget{synthesis-three-projections-of-s-ux3c6}{%
\subsection{5. Synthesis: Three Projections of (S,
φ)}\label{synthesis-three-projections-of-s-ux3c6}}

Theorem 23 establishes the stated local lattice/gauge outputs \textbf{on
the selected cubic branch} under its physical and mathematical
assumptions; uniqueness of that branch modulo an observational gauge
group is not established. M1-B and H-Bell separately condition the
existence of a measurement-independent Bell-violating completion;
Bell-inclusive uniqueness is open. The fixed-interface results of §4
supply examples of observable-preserving transformations; a general
equivalence classification and a complete projection to Standard Model
gauge symmetry are not established. This section makes the synthesis
claim explicit: quantum mechanics, general relativity, and the arrow of
time are not three independent theories but three projections of the
same finite deterministic construction.

\hypertarget{the-three-projections}{%
\subsubsection{5.1 The three projections}\label{the-three-projections}}

The framework's derivations apply the same trace-out machinery to the
same \((S, \varphi)\) along three projections: a visible-sector
projection that produces the reduced description admitting the quantum
representation ({[}Main{]}) and the Standard Model ({[}SM{]}); a
boundary-thermodynamic projection that produces general relativity
({[}GR{]}); and an arrow-of-time projection (§5.4) that produces the
emergent cosmological, thermodynamic, memory, and measurement arrows
from a single structural clock and the C1--C4 observer-selection theorem
of {[}Main, §4.6{]}.

\textbf{Visible-sector projection.} At the level of the embedded
observer's epistemic access, the trace-out over the hidden sector
produces the exact finite observable-law correspondence established in
{[}Main{]}: \(S\iff D\iff Q_{\mathrm{fb}}\). The fixed-basis
unitary/Born representation is therefore exact and bidirectional but
universal. The nontrivial OI content is history readback: C4 forces
conditional memory and hidden predictive structure in every faithful
completion; C1 and C3 diagnose the mediation and its capacity, while
C2-structural, the persistence of the record to readback, holds by
reversibility in every bijective realization and constrains
coarse-grained descriptions of the hidden sector. In the spatial theory
the dependency graph additionally supplies exact causal cones for
finite-range couplings. Full equality with the standard coherent
instrument/composite theory remains the narrower operational-lifting
statement of {[}Main §3.4{]}. Bell is separate: the
measurement-independent deterministic completion reaches quantum Bell
violation only by ontic parameter dependence, with operational
no-signaling {[}Main §3.3{]}. Applied to a cubic lattice with the wave
equation as substratum dynamics --- itself uniquely selected among
second-order reversible nearest-neighbor dynamics by center
independence, isotropy, and linearity ({[}SM, §4.1{]}) --- the selected
carrier and H-cust give the generic condensate stabilizer
\(U(3)\times U(2)\times U(1)\); the constructive central quotient and
faithful selected charged-family image give the SM Lie algebra with the
global distinctions and additional physical action premises of {[}SM,
Propositions 14a and 15a{]}, three generations (physical reading
conditional on H-spin', {[}SM §4.7{]}; chirality on H-χ', {[}SM §4.8{]})
and the Higgs as a composite scalar in the singlet staggered taste
({[}SM, §4.7{]}, Theorems 8--11), and the selected-family anomaly
classification, including its neutral-\(Q\) branch and separate
charge-normalization premise ({[}SM, §4.9{]}, Theorem 15). Twenty-two
quantitative predictions follow ({[}SM, §7{]}), including the Cabibbo
angle from a single Brillouin-zone distance, the Koide angle from a
cubic-group quadratic Casimir, all three PMNS angles within
\(1.1\sigma\) (vs NuFIT 6.0), six fermion masses from one empirical
input to better than \(1\%\), and all three SM gauge couplings at
\(M_Z\) ({[}SM, §6.3{]}).

\textbf{Boundary-thermodynamic projection.} At the level of the
partition boundary, classical horizon thermodynamics --- the
temperature, entropy, and dynamical evolution of the causal horizon as a
gravitating object --- produces general relativity through Jacobson's
thermodynamic argument, conditionally on the open covariant step ({[}GR,
§8.7{]}): the Clausius relation \(dE = T \, dS\) applied at local causal
horizons yields Einstein's equations with the same structural inputs
that give the visible-sector projection its form. Applied to the
cosmological horizon as a causal partition ({[}GR, §2{]}), this
projection determines the emergent action scale
\(\hbar = c^3 \epsilon^2 / (4G\sigma)\) from thermal self-consistency
({[}GR, §3{]}), fixes the discreteness scale
\(\epsilon = 2\sqrt{\sigma}\, l_p\), for an entropy \(\sigma\) per
boundary mode (\(\epsilon = 2 l_p\) at one nat per mode), as the unique
simultaneous solution ({[}GR, §4{]}), recovers the Bekenstein-Hawking
entropy \(S = A/(4 l_p^2)\), the mode count reproducing it once the
calibration fixes \(\epsilon/\sqrt{\sigma} = 2 l_p\) and the \(1/4\)
being the ratio \(2\pi/8\pi\) of the KMS and first-law normalizations
({[}GR, §5{]}) --- the action scale, the discreteness scale and the
coefficient conditional on H-slope with the horizon and frame conditions
of {[}GR, §2, §8.5{]}; GW250114 confirms the classical area theorem, the
coefficient itself untested directly --- and addresses the cosmological
constant problem by identifying the quantum vacuum energy and the
effective \(\Lambda\) as properties of logically distinct levels of
description rather than commensurable quantities whose \(10^{122}\)
ratio requires cancellation ({[}GR, §6{]}), the descriptive distinction
established and the dissolution conditional on the open coupling step
({[}GR, §6.2{]}). The framework carries dark energy in running-vacuum
form --- the functional form the local expansion shared by every smooth
model, the content the Type II channel with magnitude
\(|\nu| \sim 10^{-32}\) (not derived at theorem level) --- consistent
with the Bertini et al.~(2025) direct DESI-DR2-era fit at \(2\sigma\)
(which is consistent with the ΛCDM limit); a \textasciitilde{}\(95\%\)
dark sector corollary as a structural feature rather than a new particle
species; and a MOND acceleration scale \(a_0 = cH/6\) from entropy
displacement at the boundary ({[}GR, §7{]}).

The proposed projections take the same hypothesized substratum
\((S, \varphi)\) and observer partition as starting inputs. Their
distinct constructions share partial-observation machinery, but neither
the complete substratum gauge group nor all gravitational coupling steps
are derived from those inputs. They differ only in which structure of
the embedded observer's description they derive: the visible-sector
projection (§5.1, §5.3) works on the bulk dynamics and produces the
quantum-representable reduced description and the Standard Model; the
boundary-thermodynamic projection (§5.1) works on the classical thermal
data at the partition boundary and produces general relativity; the
arrow-of-time projection (§5.4) works on the horizon-clock parameter
along observer worldlines and produces the cascade of emergent arrows.

\textbf{The projections and the Poincaré structure of the partition.}
The three projections are organized by how the
observer-worldline-centered partition treats the Poincaré group, which
makes the boost/rotation result of {[}Main §3.5{]} one facet of a larger
pattern. The partition selects a preferred event --- a spatial origin
and a rest frame, jointly the observer's worldline --- so the compact
rotations about that worldline are not excluded by the partition's
rest-frame data, while the non-compact generators are broken, given that
the coarse-graining map determines the horizon foliation ({[}Main
§3.5{]}); spatial isotropy enters as an empirical input. Boosts, tied to
the selected rest frame, surface as the cosmic rest frame (operationally
the CMB frame) and the boost-sector Lorentz residual ({[}SM §3.1{]});
spatial translations, which merely relabel the observer's origin, are
broken only observer-relatively and are recovered as spatial homogeneity
--- the equivalence of all worldline-centered descriptions; and time
translation is broken intrinsically by the growth of the horizon, which
is the boundary-thermodynamic and arrow projections (§5.4, {[}GR{]}).
The partition's symmetry-breaking signature --- isotropy not broken by
the partition, a preferred rest frame, spatial homogeneity, a
distinguished cosmic time --- is the symmetry structure of a
Friedmann-Robertson-Walker cosmology. The framework does not derive that
structure; it exhibits the partition's symmetry-breaking pattern as
compatible with it, with ordinary momentum conservation holding in the
visible sector up to corrections of order (length / horizon radius).

\hypertarget{the-qm-gr-incompatibility-as-a-category-error}{%
\subsubsection{5.2 The QM-GR incompatibility as a category
error}\label{the-qm-gr-incompatibility-as-a-category-error}}

The conventional formulation of the QM-GR unification problem treats
quantum mechanics and general relativity as two competing theories in
the same logical category --- two attempts to describe the same regime
--- and asks how to merge them into a single mathematically consistent
framework. Every program of unification proceeds from this assumption:
quantum gravity programs treat the metric as a quantum field to be
quantized; geometric programs treat the wave function as a structure on
the spacetime manifold; emergent programs treat one of the two as
derived from a deeper substrate that recovers the other. The persistent
failure of all three approaches to produce a quantitatively confirmed
unification suggests that the underlying assumption is wrong.

In the present framework, the assumption is wrong because quantum
mechanics and general relativity occupy different positions in the
trace-out hierarchy. They are not two theories of the same regime at all
--- they are two projections of the same construction, viewed at
different levels. The visible-sector projection that produces the
quantum-representable reduced description and the boundary-thermodynamic
projection that produces GR are not in tension because they refer to
\emph{different objects in the construction}. The wave function is a
property of the embedded observer's compressed description of the
visible sector; the metric is a property of the classical dynamics of
the partition boundary. These are not two descriptions of the same
physical system at different levels of approximation. They are
descriptions of two different substructures of the same total object
\((S, \varphi)\).

The analogy that captures the structure most clearly is the one between
a thermodynamic and a statistical-mechanical description of a gas. The
thermodynamic description deals with pressure, temperature, and entropy;
the statistical-mechanical description deals with particle positions,
momenta, and microstate counting. These are not two competing theories
of the gas, and the question ``how do we unify thermodynamics and
statistical mechanics?'' is not a coherent question --- they describe
different things about the same system, and the relationship between
them is one of projection (statistical mechanics produces thermodynamic
quantities by averaging) rather than unification. The QM-GR relationship
in the present framework is structurally similar: GR describes the
thermodynamic (boundary, classical, deterministic) projection of
\((S, \varphi)\), and QM describes the statistical (visible-sector,
compressed, stochastic-emergent) projection. The two are related by a
definite procedure (the trace-out under conditions C1--C4), and the
question ``how do we unify them?'' is dissolved rather than answered.

This is not the same as saying that quantum mechanics is ``more
fundamental'' than general relativity, or vice versa. Both are emergent.
Both depend on the trace-out and the partition. Both are derived rather
than fundamental. The fundamental object is \((S, \varphi)\) --- a
finite set with a deterministic bijection --- and neither QM nor GR is a
feature of \((S, \varphi)\) itself. They are features of how an embedded
observer or a horizon-bounded thermodynamic regime sees \((S, \varphi)\)
from inside. (Nor is the ranking question recoverable from inside: by
the factoring argument ({[}GR §6.3{]}), no criterion grounded in the
accessible observations can separate two descriptions faithful to the
same data --- the fundamentality question is answered at the level of
\((S, \varphi)\), not adjudicated between its projections.)

\hypertarget{the-three-level-hierarchy-made-explicit}{%
\subsubsection{5.3 The three-level hierarchy made
explicit}\label{the-three-level-hierarchy-made-explicit}}

The restricted observational invariances of §4 motivate a relationship
between the visible and boundary projections; they do not establish a
complete gauge-group image under trace-out. At the substratum level the
chosen object is a bijection together with an observer partition, and
the proposed projection connects it to two separately described
structures: the emergent quantum field theory at the visible-sector
level, with the conditional generic condensate stabilizer ({[}SM,
Theorem 5{]}), the exact six-link cubic decomposition (Theorem 7), and
the constructive quotient and selected faithful charged-family image
(Propositions 14a and 15a), whose physical selection remains additional
input; and the classical horizon thermodynamics at the boundary level,
with the metric, the surface gravity, the entropy, and the temperature
determined by the partition geometry and the area-law theorem ({[}GR,
§3{]}). The exact mathematical identities in these constructions retain
their own hypotheses; a common substratum does not supply the missing
physical carrier, charged-action or gravitational coupling bridges.

The three-level structure is a proposed relationship between distinct
equivalence notions, with only the explicitly stated parts established.
Fixed-interface hidden relabellings and finite exact decoupled
extensions are visible-channel invariances (§4); the Standard Model
commutant is obtained from the cubic representation construction; and
the emergent Hamiltonian's D-gauge preserves transition data. A full
group inclusion and projection between these notions remain open, as
does the H-BG-dependent completeness of the Hamiltonian classification
{[}GR §3.3{]}.

The boundary-thermodynamic projection that produces GR is a parallel
structure at Level 2 --- but acting on the \emph{partition} rather than
on the internal field content. Where the emergent QFT is the trace-out's
effect on the bulk dynamics, the classical horizon thermodynamics is the
trace-out's effect on the partition geometry itself. Both are at Level
2. Both descend from Level 3. The framework's claim is that the SM gauge
group and Einstein's equations are not two independent theoretical
inputs but two co-derived structures, both consequences of the same
underlying \((S, \varphi)\) and the same trace-out under the same
conditions C1--C4; on the gravitational side the covariant field
equations are the open Level-G3 step ({[}GR, §8.7{]}).

\hypertarget{the-arrow-of-time-as-a-third-projection}{%
\subsubsection{5.4 The arrow of time as a third
projection}\label{the-arrow-of-time-as-a-third-projection}}

After Main §4.6 (structural observer-selection theorem), the arrow of
time joins quantum mechanics and general relativity as a third emergent
structure that the framework derives from the same underlying
\((S, \varphi)\) and the same trace-out machinery. Before §4.6, the
arrow of time was treated as external to the derivation chain --- a past
hypothesis or an anthropic selection, imposed on the framework rather
than produced by it. With §4.6 in place, the observer's confinement to
the non-equilibrium phase of \((S, \varphi)\) is a structural
consequence of the effective-mixing hypothesis, for sufficiently
fast-mixing equilibria (\(\tau_{\rm mix} \ll \tau_S\)), in the
coarse-grained description generated by iterating the one-step effective
channel, and every arrow the observer sees descends from this
confinement together with the horizon growth that the accelerating phase
produces. The third projection has the same form as the first two: a
feature of how the embedded observer sees \((S, \varphi)\) from inside,
rather than a feature of \((S, \varphi)\) itself.

A distinction should be kept in view here. The \emph{existence} of the
temporal domain --- that differentiation recurs, that there is more than
one moment at all --- is not a projection of \((S, \varphi)\) but the
framework's second foundational axiom (the recurrence axiom of the
two-axiom observation base; see {[}Main §1.2{]} and Part I of the
companion paper \emph{Physics Modulo Gauge}). What is a projection, and
what this section derives, is the \emph{arrow} --- the directedness of
time within that domain. The framework derives the arrow; it posits the
existence of the domain. Conflating the two would overstate what the
projection delivers: §5.4 produces directedness, not the bare existence
of succession.

\textbf{Horizon growth as the structural clock.} The framework's primary
arrow is the cosmological horizon entropy
\(S_{\rm dS}(t) = A(t)/(4 l_p^2)\), monotonically increasing along any
observer worldline in the \(\Lambda\)-dominated phase. The monotonicity
is a feature of the causal partition rather than of any local frame:
every observer in the accelerating phase agrees that \(S_{\rm dS}\)
increases, just as every observer in a Schwarzschild geometry agrees on
the exterior direction. The horizon clock is observer-independent,
continuous, and monotonic for as long as \(\ddot a > 0\), which the
framework identifies as the regime of interest for the derivations of
{[}GR{]}. At the substrate level, \((S, \varphi)\) has no preferred
direction: \(\varphi\) and \(\varphi^{-1}\) are equally valid as
forward-time maps. At the observer level, the horizon clock picks out a
direction, and all other arrows align with it.

\textbf{The layered cascade.} Four derived arrows align with the horizon
clock through a structural cascade. \emph{Cosmological:}
\(S_{\rm dS}(t)\) monotonic in the \(\Lambda\)-dominated phase, as
above. \emph{Thermodynamic:} via Jacobson's coupling
\(\rho_s = \rho_{\rm crit}\) ({[}GR, §3{]}), the entropy of any
horizon-bounded system increases with the observer's clock --- the
thermodynamic arrow descends from the structural arrow without an
independent postulate. \emph{Memory:} C2 (memory persistence) implies
the observer's hidden sector preserves correlations forward along the
clock rather than backward, because the coarse-grained mixing time of
the hidden-sector dynamics is uniform in the two directions but the
observer's measurement protocol commits to one. \emph{Measurement:} the
observer records ``past'' outcomes and prepares ``future'' states --- an
asymmetry inherited from the memory layer, not added on top of it. The
chain grounds in one structural fact (horizon growth) and one structural
constraint (C1--C4 observer selection via Main §4.6, conditional on its
effective-mixing hypothesis and stated for the coarse-grained
description generated by iterating the one-step effective channel). The
effective-mixing hypothesis is the premise that takes the place of an
initial-condition axiom.

\textbf{What the framework derives, and the directional-language
subtlety.} The framework's quantitative content at the level of Main
§4.6 is \emph{symmetric} in time direction --- the spectral-gap bound
\(\tau_B(V) \leq \tau_{\rm mix}(\Sigma_{\rm loc}(V)) + O(\epsilon)\),
stated for the coarse-grained description generated by iterating the
one-step effective channel, applies equally in both directions along any
observer worldline. What is \emph{imposed} by the observer, not derived,
is the convention of which direction along the horizon-clock parameter
to call ``future.'' The two directions are gauge-equivalent at the
substrate level (\(\varphi \leftrightarrow \varphi^{-1}\)), and the
observer's choice of convention is the residual freedom after all
physical content has been extracted. The framework therefore derives the
existence of a monotonic parameter along observer worldlines, not a
substrate-level arrow of time. This distinction answers Ellis's
objection (a bijective substratum cannot carry an arrow of time) by
agreeing with the objection's premise while deriving what the observer
actually sees from the structure the substratum does have.

\textbf{Scope note: causal topology as load-bearing.} The memory-arrow
step uses ``the observer's worldline has a causal topology'' as the
carrier of the asymmetry. This is not a new axiom --- {[}Main, §1.3{]}
already assumes a causal topology on the partition --- but it is
load-bearing for the cascade and is noted explicitly here rather than
treated as obvious. A framework in which the observer's worldline had no
causal topology could not run the memory-arrow step, and the cascade
would halt at the thermodynamic layer.

\textbf{Scope note: the arrow derivation is conditional on the ordering
of moments.} The cascade derives a \emph{direction} along the observer's
worldline. It presupposes that the worldline's moments are
\emph{ordered} --- that they form a sequence --- since a monotonic
parameter is monotonic with respect to an ordering. This ordering is
distinct from the bare existence of more than one moment: the
framework's foundational treatment (Part I of the companion paper
\emph{Physics Modulo Gauge}) axiomatizes the existence of a temporal
domain but treats the \emph{ordering} of its moments as a separate
question, which it leaves open. The arrow derivation of this section is
therefore conditional: it derives the direction \emph{given} an
ordering, and inherits the open status of the ordering question. What
this section establishes is the conditional --- ordered moments yield a
derived monotonic parameter and an observer-imposed direction-convention
--- not an unconditional derivation of temporal directedness from the
substratum alone. The causal topology of the worldline reduces to this
ordering together with the dynamics \(\varphi\) (it introduces no
commitment beyond the ordering), so the arrow's single open dependency
is the ordering question and nothing further.

\textbf{CP violation as the residual Ellis-class objection.} The layered
cascade closes the Ellis-axis objection for the arrow of time but does
not automatically close the related question of CP violation. The
substratum is T-symmetric (load-bearing for the reciprocity of visible
transition counts, {[}SM, §5{]}; it does not by itself fix
\(\bar\theta\)), and observed CP violation in CKM and kaon/B systems is
a flavor-basis feature of the specific bijection rather than a
consequence of the cascade. The framework treats flavor-sector CP phases
as solution-specific inputs --- see {[}SM, §5.4{]} for the scope
statement and §6.4 ``what does not falsify'' for the commitment that the
framework makes no prediction of \(\delta_{\text{CKM}}\),
\(\delta_{\text{PMNS}}\), or the Wolfenstein \((\rho, \eta)\).

\hypertarget{what-the-synthesis-does-and-does-not-claim}{%
\subsubsection{5.5 What the synthesis does and does not
claim}\label{what-the-synthesis-does-and-does-not-claim}}

The synthesis claim is structural and architectural rather than
empirical. It does not claim that the framework predicts gravitational
phenomena that conventional QFT plus GR cannot reproduce --- most of
{[}GR{]}'s predictions (the BH area law, the basic properties of horizon
thermodynamics, the broad shape of dark energy phenomenology) are also
recovered by other approaches. The framework's empirical content is
concentrated in {[}SM{]} (the SM derivation, twenty-two quantitative
predictions) and in the specific GR-side predictions of {[}GR{]} (the
Type II running-vacuum channel with \(|\nu| \sim 10^{-32}\), MOND
\(a_0 = cH/6\) with its dimensional factor imported, dark sector
concordance), which collectively distinguish the framework from standard
\(\Lambda\)CDM plus the Standard Model.

The three constructions share a candidate model and partial-observation
machinery, but use different additional premises. The cubic gauge
calculation in {[}SM{]} and the horizon mode count in {[}GR{]} are not
the same mathematical operation, and their physical interpretation is
conditional on their respective hypotheses. Their common-substratum
interpretation is a proposed relationship between the constructions, not
a proof of a unique joint physical realization.

The proposed synthesis reframes the relation between quantum
representations and gravitational thermodynamics within a candidate
deterministic model. It does not remove the physical compatibility and
reconstruction problems from the research programme. The Standard Model
gauge factors are obtained in the selected cubic representation analysis
under its own hypotheses. Identifying them as the complete image of a
full substratum gauge group requires an observational carrier and a
comparison map not established in §4. Neither that image nor uniqueness
of the physical branch follows from the common use of \((S,\varphi)\).

\begin{center}\rule{0.5\linewidth}{0.5pt}\end{center}

\hypertarget{discussion}{%
\subsection{6. Discussion}\label{discussion}}

The exact finite-law representation, conditional local construction and
scoped observational invariances address different questions. Theorem 23
organizes outputs of a selected branch; Theorem 24 distinguishes
transformations preserving coherent channels from those preserving only
populations, while all-partitions classification remains open. These
results support a structural interpretation of specified model
observables. They do not establish that mathematics exhausts reality,
that every microscopic distinction is empirically inaccessible, or that
the quantum/gravitational compatibility problem is solved.

\hypertarget{ontic-structural-realism-on-the-equivalence-class}{%
\subsubsection{6.1 Ontic structural realism on the equivalence
class}\label{ontic-structural-realism-on-the-equivalence-class}}

The structural-realist proposal of §1.1 concerns the conditional local
outputs of the selected model and invariants of a specified
observational carrier. Theorem 23 does not uniquely identify that model
among all substrata satisfying E1--E7 and A1--A6. H-Bell separately
conditions a common Bell-compatible completion; it is not the only
missing step in a full reconstruction. Realism about one complete
microscopic bijection therefore requires physical selection and
equivalence results beyond the present theorems.

This is the strongest claim the framework's theorems license at this
level without overclaiming. Two alternative framings merit explicit
comparison.

\emph{Specific-substratum realism} is the commitment that observation
singles out one microscopic bijection \(\varphi\). Neither Theorem 23's
conditional local-branch outputs nor Theorem 24's scoped invariances
establish that commitment. For a \textbf{specified} observer interface,
the proven hidden-only relabellings and exactly decoupled ancillary
enlargements preserve coherent channels; fibre-dependent hidden
relabellings preserve rooted populations without necessarily preserving
coherences. The alphabet and graph comparisons preserve only their
stated structural contents. A measurement restricted to a carrier cannot
distinguish representatives related by a transformation proved to
preserve that carrier, but no complete all-partitions gauge
classification is available to rule out further observable distinctions.

\emph{Pure explanatory economy} --- the position that the selected cubic
construction provides an economical account of several observed
structures without committing to one microscopic ontology --- remains a
defensible interpretation of the current results. Theorem 23 gives
conditional lattice/gauge outputs on that selected branch, but does not
prove that the empirical and structural inputs uniquely identify it
modulo a complete gauge group. The distinction between what the
construction explains within its assumptions and whether nature selects
that branch is substantive, rather than a question settled by
observational redundancy.

The wave function, Hilbert space, gauge fields and metric enter
particular descriptions of the candidate model, with their own observer
and physical hypotheses. This motivates comparing descriptions at
different levels, but does not prove that they are all projections of a
uniquely reconstructed observational quotient. Whether a change of
representation is physically redundant is assessed relative to the
specified observable carrier and its readout map.

Two ontological readings of the same model are indistinguishable by
measurements of a specified carrier if they stipulate identical
preparation/readout data. This is a conditional factoring statement.
Theorem 23 does not supply identical data for every candidate
microscopic interpretation, and Theorem 24 does not classify all
transformations preserving complete physical observations. A structural
reading of Wigner's puzzle is therefore an interpretation of the
construction, not an additional reconstruction theorem.

\textbf{Multi-level structural realism.} The structural-realism
commitment is not single-level. {[}Structure{]} establishes that the
framework operates at multiple levels of structural realism with
different content at each level. The hierarchy is two-dimensional: an
observation-hierarchy axis (vertical depth --- from foundational to
specific) and a gauge-hierarchy axis (horizontal breadth --- from narrow
to wide gauge equivalence). The combination produces a structure of
realisms operating at multiple intersections.

At the broadest level --- universality-class equivalence across
structural classes (Level G4 in {[}Structure §2{]}) operating at the
partial-trace observational features sub-class (Level C in {[}Structure
§2{]}) --- what is real is the algebra-channel structure of
partial-trace observation: availability of a fixed-basis Born-form
representation, channel-level unitarity in that representation, the
non-Markovian marginal in the readback sector, and the commutant
gauge-invariance pattern. These are representation/observation features
shared across the stated observer-admitting class; the claim does not
independently select the quadratic exponent or close the coherent local
operational lift.

At intermediate levels, the interpretation concerns the Standard Model
gauge algebra calculated on the selected branch and the specified
algebra/channel structures. The physical gauge-group reading retains the
carrier, amplitude-scale and matter hypotheses of {[}SM{]}; it does not
establish a unique full substratum gauge image.

At the model-specific level (Level D × Level G3 in {[}Structure §2{]}),
Theorem 23 supplies conditional outputs of the chosen cubic realization.
It does not identify every compatible model with that realization modulo
a complete gauge group. The classifications of Cabibbo, Koide and
gravitational quantities remain those of their individual derivation
chains in {[}SM{]} and {[}GR{]}. They do not establish the missing
branch selector or Bell-compatible completion.

The proposed hierarchy distinguishes model-specific assumptions from
structures preserved by specified observational comparisons. At the
local level the mathematical outputs are conditional on the selected
branch. At broader levels, identifying equivalent descriptions requires
the relevant observable structures and explicit maps. A complete
physical equivalence class and its ontological interpretation remain
separate questions.

The class-specific phenomenology is concentrated at Level D × Level G3:
the selected cubic realization, its conditional gauge/matter
interpretation and its classified quantitative comparisons. Comparisons
with other programmes must specify which observational carrier is
shared. The hierarchy supplies a way to pose these comparisons; Theorem
23 and Theorem 24 do not establish all inter-level identifications.

\hypertarget{the-ontological-hierarchy}{%
\subsubsection{6.2 The ontological
hierarchy}\label{the-ontological-hierarchy}}

The triple (S, φ, V) generates every concept in fundamental physics, not
as independent substances but as different aspects of the same
structure. \textbf{Space} is the coupling structure of φ --- the graph
G\_φ determined by which degrees of freedom affect which others ({[}SM,
§2.4{]}). \textbf{Matter} is the state --- localized patterns that
propagate through the coupling graph. \textbf{Energy} is the rate of
change under iteration. \textbf{Time} is the iteration itself.
\textbf{Quantum mechanics} is the observer's compressed description of
the visible sector. \textbf{General relativity} is the thermodynamic
limit of the coupling structure. \textbf{Conservation laws} are
emergent: energy conservation (Noether) is what information conservation
(bijectivity) looks like in the emergent quantum description. None of
these are independent entities; they are descriptions of (S, φ, V) at
different scales.

\hypertarget{the-measurement-problem}{%
\subsubsection{6.3 The measurement
problem}\label{the-measurement-problem}}

On the structural reading, the measurement problem is dissolved. The
wave function is not a component of (S, φ, V) --- it is a derived
object. Since it is derived, not fundamental, asking ``does it
collapse?'' is asking about the behavior of a compression artifact. In
the double-slit experiment, the particle traverses a single slit in the
deterministic substratum. In Wigner's friend, the Friend has a definite
outcome; Wigner's superposition reflects his epistemic deficit.

Branching is forbidden by the rigidity of φ. A fixed bijection on a
finite set has exactly one trajectory from any initial state. There is
no point at which the trajectory splits. The appearance of branching in
the emergent quantum description reflects the observer's uncertainty
about which trajectory they are on (because they cannot see the hidden
sector), not a physical splitting of worlds.

Bell correlations force a sharper distinction. For any
\textbf{finite-range Bell-relevant} substratum, the dependency graph
gives an exact causal cone; if a response-complete deterministic
completion also uses one setting-independent pre-setting ensemble, Bell
factorization follows and CHSH is capped at \(2\). The framework
therefore does not obtain quantum Bell violation from screening plus
indivisibility. It keeps measurement independence and takes ontic
parameter dependence: Bell-violating substrata have Bell-relevant
coupling range exceeding the wings' separation, while the averaged
operational statistics remain no-signaling {[}Main §3.3{]}. The
stochastic causal-local indivisible Tsirelson theorem is retained only
at that weaker stochastic layer.

\hypertarget{pre-registered-falsification-conditions}{%
\subsubsection{6.4 Pre-registered falsification
conditions}\label{pre-registered-falsification-conditions}}

The framework's structural content --- the equivalence class
\([(S, \varphi)]/\mathcal{G}_{\text{sub}}\), the emergent Standard Model
structure, and the co-derivation of quantum mechanics and general
relativity --- makes specific commitments whose future observational
status can now be stated explicitly. This subsection pre-registers the
conditions under which the framework would be falsified, separated from
properties of the specific bijection \(\varphi\) that the framework does
not claim to determine.

The distinction matters for two reasons. First, an explicitly
pre-registered falsification schedule is the empirical counterpart of
the methodological discipline developed in §§3--5: a framework that
derives rather than posits must be willing to specify what would
invalidate the derivation. Second, pre-registration separates structural
commitments --- which the framework stands or falls on --- from
solution-specific properties that lie outside its scope, and which
continued openness on does not undermine the framework.

\textbf{Class A: Structural theorems.} The following are claims about
the equivalence class itself and would falsify the framework's
reconstruction if violated.

\begin{itemize}
\tightlist
\item
  \emph{Strong CP.} The construction does not fix \(\bar\theta\) (SM
  §5.3, Theorem 21: T-invariance leaves \(\theta\in\{0,\pi\}\) and does
  not make the emergent Hamiltonian T-invariant), so the framework makes
  no commitment on the neutron electric dipole moment; reciprocity
  leaves the Jarlskog invariant free, and a mechanism making \(Y_u\) and
  \(Y_d\) simultaneously real would force \(J=0\) --- the falsification
  control on H-det proposals (SM §5.3--5.5); a measured
  \(\bar\theta\neq0\) would constrain mechanisms proposed for H-det, not
  the T-symmetry of the substratum, which is compatible with any value
  of \(\bar\theta\).
\item
  \emph{Gauge group SU(3) × SU(2) × U(1) exactly.} Any observation of a
  stable beyond-Standard-Model gauge structure --- an additional \(Z'\)
  interpretable as a new gauge boson rather than a composite resonance,
  stable exotic matter transforming under a novel gauge group, or any
  confirmed extension of the visible-sector gauge content --- falsifies
  the cubic-group decomposition of SM §4.6.
\item
  \emph{Three fermion generations.} Any observation of a stable
  fourth-generation fermion with Standard Model quantum numbers
  falsifies the taste-reduction derivation (SM §4.7). Non-stable exotic
  states consistent with composite or resonance interpretation do not.
\item
  \emph{Majorana neutrinos with normal ordering.} Inverted mass
  ordering, confirmed to high significance, falsifies the taste-breaking
  derivation of the neutrino mass matrix (SM §8.4). Confirmed Dirac
  nature of the neutrino (e.g., a null result for neutrinoless
  double-beta decay at the sensitivity that would require Majorana
  masses well below the framework's prediction) falsifies the prediction
  that no right-handed neutrino exists in the spectrum.
\item
  \emph{Bell/Tsirelson test.} For the cited \textbf{stochastic
  causal-local indivisible} sector, \(S_{\mathrm{CHSH}}\le 2\sqrt{2}\).
  At the deterministic level, measurement independence plus ontic
  parameter independence instead gives the Bell-local ceiling \(2\), so
  the framework's Bell-violating branch is ontically parameter-dependent
  {[}Main §3.3{]}. A confirmed loophole-free excess above \(2\sqrt2\)
  would falsify standard quantum mechanics and any claim that the
  operational lift lands in the standard quantum Bell set; it would not
  by itself falsify P-indivisibility. Bare C1--C4 or graph locality do
  not enforce Tsirelson.
\item
  \emph{Leading Lorentz-invariance-violation coefficient.} The
  cubic-lattice dispersion relation ({[}SM, §4.4{]} Theorem 4,
  generalized to d = 3:
  \(\cos(\omega\epsilon) = \frac{1}{3}\sum_{i=1}^3 \cos(k_i \epsilon)\)
  for a massless mode) yields standard emergent dispersion
  \(\omega = |k|\) at leading order. Series expansion gives the first
  correction as a direction-dependent quadratic term
  \(\omega^2 = k^2[1 - 2\delta(\hat k)(k\epsilon)^2 + O((k\epsilon)^4)]\),
  with \(\delta(\hat k) = \frac{1}{24}\sum_i \hat k_i^4 - \frac{1}{72}\)
  evaluated on the rigid \(\mathbb{Z}^3\) representative. The
  rigid-lattice anisotropy (values \(\delta_{[100]} = 1/36\),
  \(\delta_{[110]} = 1/144\), \(\delta_{[111]} = 0\) --- the dispersion
  is exactly linear along the body diagonals, \(\omega = k/\sqrt{3}\) to
  all orders) is a \emph{physical, direction-dependent} prediction. Its
  absolute orientation in space is gauge: generator (iv) of §4 gauges
  the rigid orientation of the cubic representative. The anisotropy
  \emph{pattern} --- its fourfold \(\ell=4\) cubic-harmonic structure
  and the direction-differences \(\delta_{[100]}-\delta_{[111]}\) --- is
  gauge-invariant, since generator (iv) preserves the octahedral
  structure and does not identify \(\mathbb{Z}^3\) with a generic
  isotropic graph (which lacks the octahedral structure the gauge group
  requires, {[}SM §4.5--4.6{]}). The orientation-averaged coefficient
  \(\langle\delta\rangle = 1/90\) is a direction-summary of this
  anisotropy, and the framework predicts the emergent dispersion
  \[\omega^2 = k^2\left[1 - \frac{(k\epsilon)^2}{45} + O\left((k\epsilon)^4\right)\right] = k^2\left[1 - \frac{4\sigma}{45}\left(\frac{E}{M_{\text{Pl}}}\right)^2 + O\left(\left(\frac{E}{M_{\text{Pl}}}\right)^4\right)\right]\]
  after using \(\epsilon = 2\sqrt{\sigma}\,l_p\) and
  \(k\epsilon = 2\sqrt{\sigma}\,E/M_{\text{Pl}}\); the coefficient of
  \((E/M_{\text{Pl}})^2\) is \(4/45\) at \(\sigma = 1\). Three
  structural features: (i) \emph{subluminal}, since
  \(\delta(\hat k) \geq 0\) everywhere on the sphere
  (\(\langle\delta\rangle > 0\) is invariant under
  orientation-averaging); (ii) \emph{no linear-in-\(E/M_{\text{Pl}}\)
  term}, because the even-parity cosines in the dispersion admit no
  \(\mathcal{O}(k^3)\) contribution --- any demonstration of linear LIV
  at any scale therefore falsifies the framework's substrate-level
  commitment to the cubic-lattice wave equation; (iii) \emph{specific
  lattice coefficient} \(1/45\) of \((k\epsilon)^2\), fixed by the
  cubic-lattice coordination and isotropy-gauge average, not tunable.
  Features (i)--(iii) and the anisotropy pattern are the Class A
  content, at tree level under the stated assumptions; the directional
  ratios
  \(\delta_{[100]} : \delta_{[110]} : \delta_{[111]} = 4 : 1 : 0\) are
  independent of \(\sigma\). The physical amplitude, the coefficient
  \(4\sigma/45\) of \((E/M_{\text{Pl}})^2\), carries the undetermined
  entropy \(\sigma\) per boundary mode and is not a parameter-free
  numerical prediction: a measurement of it would constrain \(\sigma\)
  rather than falsify the coefficient independently. Protection of the
  tree-level coefficient in the interacting theory is open ({[}SM
  §3.1{]}). Current quadratic-LIV bounds from GRB time-of-flight reach
  only \(E_{\text{QG}} \gtrsim 10^6\) GeV (roughly thirteen orders of
  magnitude below \(M_{\text{Pl}}\)), so the amplitude \(4\sigma/45\) is
  not yet near-term testable at the magnitude level; linear-LIV tests,
  by contrast, are already at sensitivities where any positive detection
  falsifies.
\end{itemize}

\textbf{Class B: Parameter-free retrodictions.} The framework produces a
set of numerical predictions that follow from the structural content
with no adjustable parameters beyond a small set of acknowledged
empirical inputs (\(m_s\) for the overall fermion mass scale, \(m_t\)
for the Higgs RGE running, \(\eta\) for the Jarlskog phase,
\(Q_{\text{down}}\) for the down-sector Koide normalization). These
predictions are stratified into five sub-classes reflecting their
structural status:

\emph{Class B-S (structural).} One prediction is a Layer 1 structural
retrodiction with no open link specific to the entry; like every Class B
entry it rests on the matter-sector hypotheses H-spin' and H-χ':

\begin{itemize}
\tightlist
\item
  The Wolfenstein parameter \(A = \sqrt{2/3}\) ({[}SM §7.1{]}).
\end{itemize}

A global-fit \(A\) outside \(\sqrt{2/3}\,(1 \pm \lambda^2)\) would
strain the leading-order identification, while the protected
\(\sin\theta_{13} = A^2\lambda\) comparison tests the geometry directly
({[}SM §7.1{]}).

\emph{Class B-C (conditional structural).} Two predictions are Layer 1
retrodictions with no fitted input whose derivation chain has one named
open link:

\begin{itemize}
\tightlist
\item
  The Cabibbo angle \(\lambda = 1/(\pi\sqrt{2})\) ({[}SM §7.1{]}) ---
  the chirality verification of the taste-changing vertex's spinor
  structure is closed via Mason et al.~(HPQCD, hep-lat/0209152),
  establishing that taste-changing transitions in staggered quarks have
  spinor structure given by a combination of \(\gamma_\mu\) and
  \(\gamma_{5\mu}\) (both chirality-preserving); the identification
  leaves the emergent-theory dressing open, and the prediction is
  conditional on it.
\item
  The Koide angle \(\theta_0 = C_2/d^2 = 2/9\) ({[}SM §7.2{]}) ---
  matched to observation (0.001\%) and the cleanest structural
  derivation (cubic-group quadratic Casimir over bandwidth squared, with
  the §7.2 uniqueness table enumerating ten alternative dimensionless
  ratios that fail to match); it is conditional with the Koide
  amplitude, whose value \(A^2 = C_2\) is not yet derived from the
  dynamics ({[}SM §7.2{]}).
\end{itemize}

A precision-upgrade violation would falsify the identification or the
closure on which the entry is conditional.

\emph{Class B-L (layered conditional, pending open derivations).} Six
predictions (the two PMNS angles counted separately) follow from the
structural form combined with one Layer-2 substrate input, one empirical
input, or one named condition on the horizon reservoir. Each has an
explicit derivational gap whose closure would move the prediction to
B-S:

\begin{itemize}
\tightlist
\item
  The mass ratio \(m_d/m_s = 1/(2\pi^2)\) ({[}SM §7.1{]}), its
  conversion from the Cabibbo scale running through the
  Gatto-Sartori-Tonin relation, an external empirical relation whose
  status the chain inherits.
\item
  The mass ratio \(m_u/m_d = \sqrt{\theta_0} = \sqrt{2/9}\) ({[}SM
  §7.2{]}) pending the explicit derivation of the ``different channels''
  mechanism producing the square root, structurally analogous to the
  open OI-vertex 1-loop derivation for \(K = 1/2\) ({[}SM §7.5{]}) in
  being mixed-layer --- substratum cubic-group structure combined with
  electroweak-symmetry-breaking emergent machinery.
\item
  The mass ratio \(m_b/m_\tau = 4.28/Z_S\) ({[}SM §7.5{]}) pending the
  explicit OI-vertex 1-loop computation confirming \(K = 1/2\), together
  with four bridge gaps in the SχPT inheritance ({[}SM §7.5{]}): the
  substratum measure reduction under Theorem 2, the 3D MC vs 4D SχPT
  dimensional bridge, the 3D BZ vs 4D Dirac taste-structure equivalence,
  and the 3D vs 4D parity identification of chirality. The per-vertex
  \(\cos^2\) symmetrization producing \(K = 1/2\) is
  dimension-independent and structurally robust; the four bridges are
  derivation gaps inherited by the other §7 mass-cluster predictions
  that share the SχPT machinery. The closure path is the two-loop SχPT
  calculation specializing Panagopoulos--Spanoudes 2017 to the OI
  cubic-group setting.
\item
  The PMNS predictions \(\sin^2\theta_{12} = 1/3 - 1/(4\pi^2)\) and
  \(\sin^2\theta_{23} = 1/2 + 1/(2\pi^2)\) ({[}SM §7.3{]}) pending
  derivation of Cond 2 (the structural relation among \(A_2\) Wilson
  coefficients) from cubic-lattice Yukawa structure. The reactor angle
  \(\sin^2\theta_{13}\) is independent of Cond 2; it inherits the
  Cabibbo \(\lambda\) and is conditional with it (class C in {[}SM
  §7.6{]}; not counted separately in Class B).
\item
  The MOND critical acceleration \(a_0 = cH/6\) ({[}GR §7.3{]}),
  inheriting the conditions on the de Sitter temperature of {[}GR
  §3.2{]} and the coupling statement \(Q = M_B c^2\) that the open G3
  map would supply, with the \(1/6\) dimensional factor imported from
  Verlinde 2016 (eq. 1.7) as \((d-3)/[(d-2)(d-1)]\) in \(d = 4\).
  Dependent predictions (baryonic Tully-Fisher, crossover radius
  \(r_M\)) inherit B-L status.
\end{itemize}

Class B-L violations at \(3\sigma\) on precision upgrade are addressed
under the pre-registration commitment via three branches: (a) closure of
the open derivation \emph{confirms} the prediction (recoverable from a
transient data fluctuation); (b) closure \emph{fails to confirm} the
prediction while preserving the framework's structural commitments
(recoverable error in the specific derivation, with the structural
commitments intact); (c) demonstration that no consistent OI-internal
derivation can produce the observed value (fatal, framework-falsifying).

\emph{Class B-M (mass chain).} Two predictions inherit a single
empirical input scale through a structural relation:

\begin{itemize}
\tightlist
\item
  The electron and muon masses \(m_e\) and \(m_\mu\) via the Koide chain
  from the empirical input \(m_\tau\) ({[}SM §7.2{]}). The prediction is
  structurally \(\theta_0 = 2/9\) with \(Q = 2/3\); the absolute scale
  is set by \(m_\tau\).
\item
  The Higgs mass \(m_H\) via SM RGE running from the structural boundary
  condition \(\lambda(M_{\text{Pl}}) = 0\) ({[}SM §7.4{]}) with \(m_t\)
  as the empirical input.
\end{itemize}

Class B-M predictions falsify under the same precision-upgrade protocol
as B-L, but the inheritance is from an empirical input scale rather than
from an open derivation.

\emph{Class B-R (retrodictions).} Two gauge-coupling values are
explicitly retrodictions per the §6.3 {[}SM{]} parameter count:

\begin{itemize}
\tightlist
\item
  The SU(2) gauge coupling at \(M_Z\) ({[}SM §6, items 11-12 in
  §7.6{]}).
\item
  The SU(3) gauge coupling at \(M_Z\).
\end{itemize}

Three fitted parameters \((\delta_0, A, B)\) against three observed
couplings produces zero residual by construction; these classifications
are honest. Class B-R predictions can never falsify in the
precision-upgrade sense --- their value is what falsification of the
\emph{framework's gauge-sector structural commitments} (universality of
\(\delta_0\), the resummation form, the structural prediction
\(1/\alpha_0 = 23.25\) from the 1-loop staggered VP (ambient
coefficient; native for hypercharge only under H-observer-bundle and
H-Y-vertex, {[}SM §6.5{]}), and the \(A \cdot B\) cross-check
(\(\approx 48\) under the \(1/N^2\) leading-power assumption;
inconclusive across finite-size models, {[}SM §6.2.1{]})) would test.

\emph{Status summary.} Of the predictions in Class B: \textbf{1 is
structural with no open link of its own} (B-S), \textbf{2 are
conditional structural} (B-C), \textbf{6 are layered conditional pending
open derivations} (B-L), \textbf{2 are mass-chain inheritances} (B-M),
and \textbf{2 are retrodictions} (B-R). Total 13 predictions in Class B.
Their empirical standing is assessed entry by entry: the gauge couplings
match by construction, and the quoted precision of a measured comparison
does not apply to every entry.

The current status of each is documented in the SM and GR companion
papers with experimental references. As experimental precision improves,
the discrimination power of the B-S set increases monotonically; closure
of the open derivations would move B-C and B-L predictions to B-S over
time.

The Bekenstein-Hawking coefficient \(S = A/(4 l_p^2)\), with the factor
\(1/4\) the ratio \(2\pi/(8\pi)\) between the Euclidean KMS period and
the Jacobson Einstein-equation prefactor ({[}GR §5{]}), is a consistency
relation of the horizon calibration and is not counted among them. It is
conditional on the infrared detailed-balance condition H-slope on the
horizon reservoir together with the horizon and frame conditions of
{[}GR §2, §8.5{]}; the counting \(S = \sigma A/\epsilon^2\), for an
entropy \(\sigma\) per boundary mode, is unconditional and the
coefficient enters through \(\epsilon^2 = 4\sigma\,l_p^2\), from which
\(\sigma\) cancels. The closure path is the continuum-generator theorem
of {[}GR §2{]}. GW250114 confirms the classical area theorem, which the
framework preserves; the \(1/4\) coefficient has no direct empirical
test.

\textbf{Class C: Framework-level commitments.} Two broader commitments
are implicit in the construction; each is stated with what would
contradict it.

\begin{itemize}
\tightlist
\item
  \emph{Finite classical realizability of quantum processes.} The
  characterization theorem ({[}Main, §3.4{]}) establishes that the
  quantum representation is the description admitted by an embedded
  observer satisfying C1--C4 on a deterministic substrate with bounded
  classical hidden-sector memory. At finite accuracy, every finite
  adaptive quantum experiment has a finite deterministic realization
  (the gluing and uniform-continuum theorems of {[}Main, §3.4{]}), while
  the exact continuum theory at any fixed finite carrier is excluded by
  the classical-dimension obstruction. Non-Markovian models are
  permitted and expected (they are how the hidden-sector correlation
  persistence manifests). No finite-accuracy experiment contradicts this
  commitment by itself. A memory-witness experiment would contradict the
  ontology only given the coherent-preparation and instrument bridge, an
  outcome-independent rule identifying fundamental open-system
  processes, and a constraint derived from the completed model on what
  those processes must show; the last is not established ({[}Main,
  §4.5{]}).
\item
  \emph{The dark sector as entropy-displacement, not as a novel particle
  content.} The framework's account of dark energy as boundary-entropy
  bookkeeping read in the Type II RVM channel, and dark matter as frozen
  boundary entropy (GR §7.3), is incompatible with the direct detection
  of a WIMP interpretable as a fundamental particle with
  Standard-Model-like interactions and no entropy-displacement
  signature. A confirmed WIMP with such properties falsifies the
  framework's dark-sector account while leaving the rest of the
  structure intact --- it would require abandoning the §7.3 derivation
  and either identifying an error in it or accepting that the
  framework's dark-sector story is wrong even as other derivations
  survive.
\end{itemize}

\textbf{Cosmological commitment.} The dark energy sector is subtler than
Classes A--C. The framework supplies geometry-dependent dark-energy
bookkeeping whose smooth late-time expansion contains an \(H^2 - H_0^2\)
term --- the local form shared by every smooth model --- and reads it in
the Type II running-vacuum channel; it does not presently derive a
distinctive Type II law (GR §7.1). The magnitude is given by the
emergent-QFT calculation
\(|\nu_{\rm QFT}| \sim (M_{\rm SM}/M_{\rm Pl})^2 \sim 10^{-32}\)
(conditional on the GR §7.3 split, not yet established at theorem
level), observationally indistinguishable from zero. Substratum
contributions beyond this are constrained by the boundary mode count of
§3-§5, which catalogues modes by spatial location alone (one mode per
\(\epsilon^2\) cell, no inherent frequency assignment --- a count forced
by the link-direction structure of the \(K = 2d\) components, {[}GR
§3.1{]}); the natural 2D distribution concentrates them near the UV
cutoff, contributing well below \(|\nu_{\rm QFT}|\). The framework's
commitment is therefore the Type II reading, with magnitude
\(|\nu| \approx |\nu_{\rm QFT}|\) (conditional, per GR §7.1), and the
Bertini et al.~(2025) measurement
\(\nu = -(2.5 \pm 1.3) \times 10^{-4}\) consistent with this expectation
at \(2\sigma\), where Bertini is itself consistent with zero. The
falsification band is correspondingly stratified: a DESI Year 5
measurement of \(|\nu| \gtrsim 10^{-4}\) at high significance would
require additional substratum structure (log-uniform mode-per-decade
distribution at the cosmological horizon) not currently in §3-§5; a
measurement of \(\nu\) with high-significance positive sign would
falsify the Type II functional form itself, which is a structural
commitment rather than a magnitude one. The first outcome is a
refinement within the current bracket; the second is a structural
falsification.

\textbf{Unfixed inputs and falsification.} The absolute mass scales,
unfixed CP phases, initial conditions and other explicitly supplied
model inputs are not parameter-free predictions. Once a model fixes
them, it must still satisfy the empirical constraints and any structural
relations claimed for that model. Reciprocity of transition counts fixes
neither \(\bar\theta\) nor the Jarlskog invariant {[}SM §5.3--§5.4{]}.
The dynamics \(\varphi\) is not a microstate: changing it can change
observable predictions. Theorem 24 proves specified transformations
observationally invisible on specified carriers; it does not classify
every choice of \(\varphi\) as gauge. A hidden initial state can also
affect later observations. Only differences proved equivalent for the
chosen preparation/readout carrier are empirically redundant relative to
that carrier.

\textbf{The pre-registration commitment.} The author commits to treating
the Class A conditions, the Class C commitments, and a Class B
retrodiction exceeding \(3\sigma\) on precision upgrade as
framework-falsifying rather than as an invitation to post-hoc rescue. In
the case of a Class B violation, the framework will be revisited to
determine whether the violation reflects a computational error in the
specific retrodiction, an incorrect specific-bijection input, or a
structural problem with the derivation; the first two are recoverable,
the third is fatal. The framework does not reserve the right to add
structural postulates after the fact to accommodate disconfirming
evidence. Pre-registration of this commitment is itself part of the
falsifiability content of the framework.

\textbf{Regimes of deliberate silence.} Three regimes lie outside the
framework's predictive scope by construction and are not objects of
falsification claims. First, the \emph{pre-horizon early universe} ---
the regime before a stable causal partition with \(\tau_B \sim H^{-1}\)
exists --- has no observer in the framework's sense ({[}Main, §1.3{]})
and no horizon-level structure for the {[}GR{]} derivations to act on;
the framework's only commitment in this regime is a boundary condition
on deeper theories, namely that any successful pre-horizon dynamics must
produce a C1--C4-satisfying regime by the time a horizon stabilizes.
Second, the \emph{equilibrium phase} of \((S, \varphi)\) is, by the
structural observer-selection theorem, observer-free for sufficiently
fast-mixing equilibria (\(\tau_{\rm mix} \ll \tau_S\)) in the
coarse-grained description generated by iterating the one-step effective
channel, conditional on its effective-mixing hypothesis ({[}Main,
§4.6{]}); the framework makes no prediction for this regime because it
has no observers to whom a prediction could be addressed. Third, regions
\emph{beyond the observer's causal horizon} are part of that observer's
hidden sector and are addressed only through boundary effects (the
area-law entropy of {[}GR, §5{]}); the framework takes no
observer-independent position on their internal content. These silences
are structural rather than methodological --- they follow from the
characterization theorem's identification of QM with embedded
observation, and from the observer-selection theorem's restriction of
observers to the non-equilibrium phase --- and stating them explicitly
sharpens the scope of the positive commitments above.

\begin{center}\rule{0.5\linewidth}{0.5pt}\end{center}

\hypertarget{conclusion}{%
\subsection{7. Conclusion}\label{conclusion}}

Theorem 23 collects conditional local-branch consequences under the
selected wave-rule, H-link, representation, amplitude-scale and matter
premises, with the gravitational calibration retaining its horizon
conditions. It does not prove the converse from E1--E7/A1--A6 to a
unique local substratum. Theorem 24 establishes exact hidden-only
channel invariance, broader rooted-population invariance and
well-defined independent ancillary enlargement, while its alphabet/graph
comparisons retain their structural scope. Pair C refutes exhaustiveness
of the restricted coherent-channel moves. Proposition 24.2 separately
gives fixed-carrier hidden-unitary rigidity from all block-word traces
and the spanning premise; visible channel equality alone does not supply
those hypotheses. The all-partitions classification, the full Standard
Model gauge-image map and a common Bell-compatible physical completion
remain open. The quantum, gauge and horizon constructions motivate a
common-substratum interpretation, without establishing a unique full
observational quotient.

The synthesis claim that follows is not a unification of quantum
mechanics and general relativity in the traditional sense. The
traditional unification problem treats the two theories as competing
accounts of the same regime and asks how to merge them. The present
framework treats them as projections of the same construction onto two
different substructures --- the visible-sector trace-out and the
boundary-thermodynamic limit --- and the question of how to merge them
is dissolved as a category error rather than answered as a technical
problem. The result is not a theory of everything but a structural
account of why the apparent incompatibility between quantum mechanics
and general relativity has been so resistant to resolution: there is
nothing to resolve, because the two are not in competition. They are
co-derived from the same object by the same procedure under the same
conditions, and the framework developed across this four-paper synthesis
({[}Main{]}, {[}SM{]}, {[}GR{]}, {[}Substratum{]}) makes this exact
rather than approximate.

Three sets of open problems remain. Neither these results nor this
four-paper synthesis require the framework to be the final theory of
physics: the substratum may itself be an effective description of
something deeper, and the trace-out may have corrections beyond the
leading-order results developed here. First, the \emph{specific
Bell-inclusive completion} that describes our universe is not uniquely
determined by the framework. Theorem 23 gives conditional local
lattice/gauge outputs; selection of a unique branch modulo a full gauge
classification and the preparation-indexed Bell-composite rule remain
open. The framework predicts the structural features of the Standard
Model and the gravitational sector but not the specific values of (for
example) the lightest fermion mass or the specific numerical value of
the Type II RVM coefficient \(\nu\) beyond its dimensional order of
magnitude and sign. These are properties of the particular bijection,
analogous to the mass of the sun in general relativity. Second, the
\emph{trans-Planckian regime} --- the regime in which the lattice
spacing \(\epsilon = 2\sqrt{\sigma}\,l_p\) is not small compared to the
relevant length scale --- lies outside the framework's leading-order
results. The framework presents the lattice as fundamental, not as a
regulator approximating a continuum theory, and so does not need a
continuum limit in the standard sense; but the corrections to the
leading-order results in the regime where the partition geometry varies
on lattice scales are an open question. Third, the \emph{initial
conditions} --- the specific configuration of the bijection at any given
time --- are likewise not determined by the framework. The framework
establishes which structures emerge from the bijection but not which
microstate the universe is currently in.

These open problems are framed within the construction rather than
against it. Each is sharply formulated and admits a definite (if
presently unanswered) form. The progress reported here is that the
structural relationship between quantum mechanics and general relativity
--- the central open problem of fundamental physics for nearly a century
--- is fixed by the construction developed across this four-paper
synthesis, and the remaining open problems are problems within the
construction rather than problems with it.

\begin{center}\rule{0.5\linewidth}{0.5pt}\end{center}

\hypertarget{references}{%
\subsection{References}\label{references}}

{[}1{]} R. Bousso, ``The holographic principle,'' \emph{Rev.~Mod. Phys.}
\textbf{74}, 825 (2002).

{[}2{]} N. Bao, S. M. Carroll, and A. Singh, ``The Hilbert space of
quantum gravity is locally finite-dimensional,'' \emph{Int. J. Mod.
Phys. D} \textbf{26}, 1743013 (2017).

{[}3{]} J. Maldacena and L. Susskind, ``Cool horizons for entangled
black holes,'' \emph{Fortsch. Phys.} \textbf{61}, 781--811 (2013);
arXiv:1306.0533.

{[}4{]} S. Pedalino, B. E. Ramírez-Galindo, R. Ferstl, K. Hornberger, M.
Arndt, and S. Gerlich, ``Probing quantum mechanics with nanoparticle
matter-wave interferometry,'' \emph{Nature} \textbf{649}, 866 (2026).

{[}5{]} J. Berstel and C. Reutenauer, \emph{Noncommutative Rational
Series with Applications}, Cambridge University Press, Chapter 2,
``Minimization,'' Theorem 2.4; doi:10.1017/CBO9780511760860.003.

{[}6{]} M. Thon and H. Jaeger, ``Links between multiplicity automata,
observable operator models and predictive state representations: a
unified learning framework,'' \emph{Journal of Machine Learning
Research} \textbf{16}, 103--147 (2015).

{[}7{]} W. Specht, ``Zur Theorie der Matrizen II,'' \emph{Jahresbericht
der Deutschen Mathematiker-Vereinigung} \textbf{50}, 19--23 (1940).

{[}8{]} N. A. Wiegmann, ``Necessary and sufficient conditions for
unitary similarity,'' \emph{J. Austral. Math. Soc.} \textbf{2}, 122--126
(1961), Thm. 4; doi:10.1017/S1446788700026422.

{[}9{]} V. Futorny, R. A. Horn and V. V. Sergeichuk, ``Specht's
criterion for systems of linear mappings,'' \emph{Linear Algebra and its
Applications} \textbf{519}, 278--295 (2017), §1 and Corollary 2;
doi:10.1016/j.laa.2017.01.006; arXiv:1701.08826.

{[}10{]} M. P. Schützenberger, ``On the definition of a family of
automata,'' \emph{Information and Control} \textbf{4}, 245--270 (1961);
doi:10.1016/S0019-9958(61)80020-X.

{[}11{]} J. W. Carlyle and A. Paz, ``Realizations by stochastic finite
automata,'' \emph{Journal of Computer and System Sciences} \textbf{5},
26--40 (1971); doi:10.1016/S0022-0000(71)80005-3.

{[}12{]} M. Fliess, ``Matrices de Hankel,'' \emph{Journal de
Mathématiques Pures et Appliquées} \textbf{53}, 197--222 (1974); erratum
\textbf{54}, 481 (1975).

\begin{center}\rule{0.5\linewidth}{0.5pt}\end{center}

\textbf{Companion papers (cited inline by short name):}

{[}Main{]} A. Maybaum, ``The Incompleteness of Observation,'' (2026).

{[}SM{]} A. Maybaum, ``The Standard Model from a Cubic Lattice,''
(2026).

{[}GR{]} A. Maybaum, ``ℏ, the Bekenstein-Hawking Entropy, and the
Running-Vacuum Form of Dark Energy from the Cosmological Horizon,''
(2026).

\end{document}
% source-sha256: a8daee0844982283d99eef518111a2944b19046ba8b5c1a0f475381981e2f651
